% !TeX root = ../../Book3.tex
% 纲要标题：## 第22章*　局部上同调、典范模与 Gorenstein 环
% 纲要位置：计划纲要.md，第 2667 行。

\OptionalChapter{局部上同调、典范模与 Gorenstein 环}{b3:ch:22}
\BookChapterNumber{b3:ch:22}

\begin{chapterabstract}
本章由 Čech 复形计算局部上同调，并用首次非零次数刻画深度。完备局部环上的 Matlis 对偶，将 Artin 模与有限生成模相互联系。把局部环写成完备正则局部环的商后，环境环上的 Ext 给出典范模；Cohen–Macaulay 条件使对偶信息集中在一个次数，Gorenstein 自对偶\footnote{戈伦斯坦（Daniel Gorenstein，1923-1992），美国数学家，研究有限单群分类；其早年关于平面曲线的工作引出了以他命名的环。}则进一步表现为典范模与环自身同构。非 Cohen–Macaulay 情形须保留多个对偶次数。
\end{chapterabstract}

\begin{learninggoals}
\begin{itemize}
\item 用有限 Čech 复形计算局部上同调，并解释深度和维数给出的两端消失区间。
\item 检验内射包络的本质扩张条件，写出有限生成模与 Artin 模的 Matlis 求值同构，并指出完备性的作用。
\item 用来自完备正则局部环的满射表示构造典范模，解释 CM 条件如何使局部对偶集中在一个次数。
\item 由 Koszul 对偶计算完全交的典范模，用基座或具体乘法配对检验零维 Gorenstein 性。
\end{itemize}
\end{learninggoals}

设 \(k\) 为域，令 \(R=k[x,y]\)，把 \(R_x,R_y\) 放在同一个 Laurent 多项式环 \(R_{xy}\) 中。\(1/(xy)\) 不能写成这两个子环中元素之差：\(R_x\) 的单项式有非负的 \(y\) 指数，\(R_y\) 的单项式有非负的 \(x\) 指数，故两者之和中 \(x^{-1}y^{-1}\) 的系数总为零。这就给出 \(R_{xy}/(R_x+R_y)\) 的一个非零类，显示交叠处的分式未必来自两边局部化数据的差。

\ProseParagraphBreak

Čech 复形把这种“局部有解、全局不能拼合”的差异延伸到多个开集，并产生局部上同调。它的首次非零次数与前面用正则序列定义的深度相遇；而在完备局部情形，与内射模中的对偶计算相配，又可把上同调转成 Ext 的信息。这些 Ext 模是有限生成的，因而原本含有无界分母的局部上同调，可以借助有限生成模的消解来研究。

\ProseParagraphBreak

本章需要\chapref{b3:ch:13}的 Artin–Rees 与完备化、\chapref{b3:ch:15}的 Koszul 复形、\chapref{b3:ch:16}的深度，以及\chapref{b3:ch:18}和\chapref{b3:ch:19}的正则局部环与 CM 环。第二卷第7.2节的 Baer 判据、足够多内射模定理也将直接使用。内射消解、比较映射及单行双复形计算在使用处给出构造；一般导出函子和导出范畴的系统论述安排在第四卷。

\ProseParagraphBreak
设 \(R\) 是来自完备正则局部环 \(S\) 的满射像。典范模可以从 \(\operatorname{Ext}_S(R,S)\) 的一个次数取出；若 \(R\) 为 Cohen–Macaulay 环，其余次数消失，便得到只涉及这个模的局部对偶。完全交的 Koszul 消解具有自对偶性，因此典范模同构于 \(R\)；不满足 CM 条件的商环则可能保留多个非零次数。完备局部环具有这类满射表示，是 Cohen 结构定理的结论，见\appref{b3:app:D}。

\ProseParagraphBreak
局部上同调的导出函子定义、Koszul 计算及深度判据，可参阅 Weibel《An Introduction to Homological Algebra》\cite[\S 4.6, pp. 115--119]{Weibel1994HomologicalAlgebra}。一般 Gorenstein 环的判别见松村英之《Commutative Ring Theory》\cite[\S 18, Theorem 18.1, pp. 141--144]{Matsumura1989CommutativeRingTheory}；内射包络和 Matlis 对偶见同书\cite[\S 18, Theorems 18.4--18.6, pp. 145--149]{Matsumura1989CommutativeRingTheory}。典范模与局部对偶的进一步论述，见 Bruns 与 Herzog《Cohen--Macaulay Rings》修订版\cite[\S\S 3.3, 3.5]{BrunsHerzog1998CohenMacaulay}。

% !TeX root = ../../Book3.tex
% 纲要标题：### 22.1 局部上同调
% 纲要位置：计划纲要.md，第 2669 行。

\section{局部上同调}
\label{b3:sec:2201}

模中的元素可能在局部化后消失，几份局部化中的元素也可能无法同时由原模中的元素给出。支集挠子函子的导出把这些现象按次数区分，有限 Čech 复形则给出可直接计算的表示。


% 纲要内容（仅作写作计划，尚非教材正文）：
%
% * 支集挠子函子（即本章的 $I$-挠子函子 $\Gamma_I$，英文 torsion functor；不是 torsor）
% * 内射消解负责导出定义，有限 Čech 复形负责计算；双复形两种计算、二元算例与平坦基变换
% * 逐元素幂挠与统一幂零化的区别；正则元商的深度归纳与参数个数给出的高次消失
%
% ---
%

\subsection{支集挠子函子}
\label{b3:subsec:2201:01}

局部化 \(M_f\) 保留在 \(f\ne0\) 的地方能够看见的元素；映射 \(M\to M_f\) 的核则记录离开 \(V(f)\) 后消失的部分。对于一个由多个元素生成的理想，除了这种核，还会出现局部数据不能拼合的障碍。局部上同调把这些障碍按次数组织起来。

\begin{definition}\label{b3:def:support-torsion-functor}
设 \(R\) 为交换 Noether 环，\(I\subseteq R\) 为理想，\(M\) 为 \(R\) 模。令
\[
\Gamma_I(M)\coloneqq\{m\in M:\text{存在}\;n\ge1\;\text{使}\;I^nm=0\}.
\]
把 \(M\) 送到这个子模、把模同态送到其限制，得到的函子称为\term[zhijinaozihanzi]{支集挠子函子}[torsion functor]，也称 \(I\)-挠子函子。
\end{definition}
\symidx{\(\Gamma_I(M)\)：支集挠子模}

这里的“支集”有直接含义：\(m\in\Gamma_I(M)\) 等价于循环模 \(Rm\) 的支集包含于 \(V(I)\)。事实上，循环模的支集为 \(V(\operatorname{Ann}_R(m))\)，包含关系等价于 \(I\subseteq\sqrt{\operatorname{Ann}_R(m)}\)；有限组生成元各有一个幂落入零化子，取足够大的总次数就得到 \(I^nm=0\)。定义中的量词是“对每个元素存在一个指数”，这个指数可以随元素改变；即使 \(\Gamma_I(M)=M\)，也未必存在一个统一的 \(N\) 使 \(I^NM=0\)。

\ProseParagraphBreak
若 \(0\to M'\to M\to M''\) 正合，则 \(0\to\Gamma_I(M')\to\Gamma_I(M)\to\Gamma_I(M'')\) 正合：前两处只是在原模中检验相同的零化条件。因此它是左正合函子；但右端未必满射。例如对域 \(k\)，在 \(R=k[x]\)、\(I=(x)\) 下，满射 \(R\to R/(x)\) 的两端挠子模分别为零与 \(k\)。为记录这一失败，需要把原函子延伸到正次数。

\ProseParagraphBreak
内射消解使这一延伸成为 \(\Gamma_I\) 的右导出函子，并由比较映射和短正合列的消解得到函子性与长正合列。随后要证明的 Čech 比较定理承担另一项任务：把这个由内射消解定义的对象写成有限的局部化复形，供具体计算使用。

\begin{definition}\label{b3:def:local-cohomology}
设 \(R\) 为交换 Noether 环，\(I\subseteq R\) 为理想，\(M\) 为 \(R\) 模。选一个内射消解
\[
0\longrightarrow M\longrightarrow E^0\xrightarrow{\mathrm{d}^0}E^1
\xrightarrow{\mathrm{d}^1}E^2\longrightarrow\cdots.
\]
定义 \(M\) 关于 \(I\) 的第 \(i\) 个\term[jubushangtongdiao]{局部上同调模}[local cohomology module]为
\[
H_I^i(M)\coloneqq H^i(\Gamma_I(E^\bullet))\quad(i\ge0),
\]
并令负次数的局部上同调为零。亦记 \(H_I^i=R^i\Gamma_I\)，即 \(\Gamma_I\) 的第 \(i\) 个右导出函子。
\end{definition}
\symidx{\(H_I^i(M)\)：局部上同调模}

第二卷第7.2节的足够多内射模定理已保证这种消解存在：先嵌入 \(E^0\)，再把余核嵌入 \(E^1\)，逐次进行。消解之间的比较也由目标项的内射性构造。给定 \(M\to N\)，先将它延拓到次数零；在下一次数，先在前一微分的像上规定映射，再延拓到整个项。链映射关系保证这个规定与核相容，因而可逐次完成。

\ProseParagraphBreak
若两条比较链映射提升同一个 \(M\to N\)，它们的差在次数零的 \(M\) 上为零，故先经过 \(E^0/M\cong\operatorname{im}(E^0\to E^1)\) 分解，再利用目标次数零项的内射性将该映射延拓到 \(E^1\)，得到第一项同伦。逐次减去已构造的同伦项，在下一次数重复，得到差为 \(\mathrm{d}h+h\mathrm{d}\)。加性函子 \(\Gamma_I\) 将这个同伦限制到挠子子模，所以两条链映射诱导相同的上同调映射。将此用于恒等映射的两次延拓，便得到不同消解之间互逆的同构；用于一般模同态，则得到所述函子性。

\begin{proposition}\label{b3:prop:local-cohomology-basic}
设 \(R\) 为交换 Noether 环，\(I,J\subseteq R\) 为理想，\(M\) 为 \(R\) 模。则 \(H_I^0(M)=\Gamma_I(M)\)。短正合的 \(R\) 模列给出各阶 \(H_I^i\) 的长正合列；若 \(\sqrt I=\sqrt J\)，则 \(H_I^i(M)\cong H_J^i(M)\) 自然成立。
\end{proposition}
\begin{proof}
零次结论就是左正合性。对于 \(0\to M'\to M\to M''\to0\)，给两端选内射嵌入 \(M'\to E'\)、\(M''\to E''\)。把第一条映射延拓到 \(M\)，再合上 \(M\to M''\to E''\)，得到单射 \(M\to E'\oplus E''\)。其余核与两端的余核成短正合列，故可重复此过程，得到逐项分裂的三份内射消解。加性函子 \(\Gamma_I\) 保持这些逐项分裂列，圈的提升、取微分和取商遂给出长正合列，连接同态的良定义及正合性与\cref{b3:lem:ext-basic-properties}中的复形论证相同。

若两理想根相等，有限生成性给出 \(I^a\subseteq J\)、\(J^b\subseteq I\)。因此两种挠子条件相同，\(\Gamma_I=\Gamma_J\)；对同一内射消解取上同调即可。
\end{proof}

\subsection{Čech 复形与导出函子的比较}
\label{b3:subsec:2201:02}

内射消解适合定义，但往往不便计算。对于 \(I=(f_1,\ldots,f_r)\)，仅用局部化就可以写出一个有限复形。

\begin{definition}\label{b3:def:cech-complex}
设 \(R\) 为交换环，\(f_1,\ldots,f_r\in R\)，\(M\) 为 \(R\) 模。对任意 \(f\in R\)，取两项上链复形 \(C^\bullet(f;R)=[R\to R_f]\)，次数为零和一。令
\[
C^\bullet(f_1,\ldots,f_r;M)
=M\otimes_R C^\bullet(f_1;R)\otimes_R\cdots\otimes_R C^\bullet(f_r;R),
\]
张量微分在越过次数 \(q\) 的因子时带符号 \((-1)^q\)。这个复形称为所给元素组关于 \(M\) 的\term[Cechfuxing]{Čech 复形}[Čech complex]。
\end{definition}
\symidx{\(C^\bullet(f_1,\ldots,f_r;M)\)：Čech 复形}

具体地，第 \(p\) 项为所有 \(M_{f_{i_1}\cdots f_{i_p}}\) 的有限直和；次数零为 \(M\)。微分是带交错符号的局部化映射。两元情形为
\[
0\longrightarrow M\xrightarrow{m\mapsto(m,m)}M_f\oplus M_g
\xrightarrow{(a,b)\mapsto b-a}M_{fg}\longrightarrow0.
\]
这里零次上同调是那些在两份局部化中都消失的原元素。一次上同调由在 \(M_{fg}\) 中相等的元素对组成，再商去来自 \(M\) 的对角像，因而测量相容的局部数据能否由同一个原元素给出。二次上同调则是 \(M_{fg}/(\operatorname{im}M_f+\operatorname{im}M_g)\)，测量交叠处的元素能否写成两份局部数据之差。若以空元素组生成零理想，则复形只有次数零的 \(M\)。有限复形并不意味着其上同调有限生成：\(k[x,x^{-1}]/k[x]\) 已经是一个反例。

\ProseParagraphBreak
同一组局部化还会在\secref{b3:sec:2405}的开覆盖计算中出现，但起始次数不同。对 \(I=(f,g)\)，两种排列为
\[
\begin{array}{c|ccccc}
\text{次数}&0&&1&&2\\\hline
\text{支集复形}&M&\longrightarrow&M_f\oplus M_g&\longrightarrow&M_{fg}\\
\text{开覆盖复形}&M_f\oplus M_g&\longrightarrow&M_{fg}&\longrightarrow&0
\end{array}
\]
各行的箭头都是上面的局部化映射。第二行删去第一行的 \(M\)，再将其余项左移一位，因此正次数的上同调与第一行相差一次。这里的两行都只是局部化模构成的复形；第二行作为 \(D(f)\cup D(g)\) 开覆盖模型的解释，以及它的零次项与一般开集截面的关系，将在\secref{b3:sec:2405}定义相关几何对象后写成正合列。

\begin{lemma}\label{b3:lem:injective-cech-acyclic}
设 \(R\) 为交换 Noether 环，\(I\subseteq R\) 为理想，\(E\) 为内射 \(R\) 模。则 \(\Gamma_I(E)\) 仍内射。对任意 \(f\in R\)，映射 \(E\to E_f\) 满射。因此，任意生成组 \(I=(f_1,\ldots,f_r)\) 满足
\[
H^0(C^\bullet(f_1,\ldots,f_r;E))=\Gamma_I(E),\qquad
H^i(C^\bullet(f_1,\ldots,f_r;E))=0\quad(i>0).
\]
\end{lemma}
\begin{proof}
先证明挠子部分仍然内射，再证明局部化映射满射。这使一元 Čech 复形在内射模上只留下零次上同调；由于挠子部分仍内射，多元情形便能逐个加入其余平坦的局部化因子。

用第二卷第7.2节已经证明的 Baer 判据检验第一项。给理想 \(\mathfrak a\subseteq R\) 和同态 \(u:\mathfrak a\to\Gamma_I(E)\)。因 \(\mathfrak a\) 有限生成，某个 \(I^n\) 零化全部像。对有限生成的子模 \(\mathfrak a\subseteq R\) 应用\cref{b3:thm:artin-rees}，给出足够大的 \(N\)，使
\(\mathfrak a\cap I^N\subseteq I^n\mathfrak a\)。Artin–Rees 在这里把“像被 \(I^n\) 零化”转成一个统一的交理想条件，使延拓能够在 \(I^N\) 上取零。于是 \(u\) 在该交上为零，可延拓为
\(\mathfrak a+I^N\to E\)、\(a+b\mapsto u(a)\)。再用 \(E\) 的内射性延拓到 \(R\)。设一的像为 \(e\)，则 \(I^Ne=0\)，所以这个延拓实际取值于 \(\Gamma_I(E)\)。

为了证明满射性，给 \(e/f^n\in E_f\)。环的升链条件使 \(0:_Rf^k=0:_Rf^{k+n}\) 对充分大的 \(k\) 成立。因而
\[
f^{n+k}R\longrightarrow E,\qquad f^{n+k}a\longmapsto f^kae
\]
良定义。延拓到 \(R\to E\)，设一的像为 \(z\)，则 \(f^{n+k}z=f^ke\)，即 \(z/1=e/f^n\)。当 \(f\) 幂零时局部化为零，也包含在此证明中。

现在对 \(r\) 归纳。已证的一元情形给出短正合列
\[
0\longrightarrow\Gamma_{(f_1)}(E)\longrightarrow E\longrightarrow E_{f_1}\longrightarrow0.
\]
因而次数零的 \(\Gamma_{(f_1)}(E)\) 嵌入两项复形 \([E\to E_{f_1}]\)，商复形正合。将这一复形短正合列与余下 Čech 复形的每个平坦项张量，得到逐项短正合列，且商双复形的各行仍正合，故总复形的上同调等于
\(C^\bullet(f_2,\ldots,f_r;\Gamma_{(f_1)}(E))\) 的上同调。这里利用这些正合行逐次消去圈的分量，有限项数保证终止；所用微分符号正是张量符号。第一项已经证明中间模内射，归纳假设可用。最后，先被 \(f_1\) 的某幂零化、再被其余生成元各自的某幂零化，等价于被 \(I\) 的某幂零化，所以零次为 \(\Gamma_I(E)\)。
\end{proof}

\begin{theorem}\label{b3:thm:cech-local-cohomology}
设 \(R\) 为交换 Noether 环，\(I=(f_1,\ldots,f_r)\)，\(M\) 为任意 \(R\) 模。则有自然同构
\[
H_I^i(M)\cong H^i(C^\bullet(f_1,\ldots,f_r;M))\qquad(i\ge0).
\]
特别地，右侧不依赖所选生成组，每个 \(H_I^i(M)\) 都是 \(I\)-挠子模。
若 \(R\to B\) 为交换 Noether 环之间的平坦同态，则对所有 \(i\ge0\) 有自然同构
\[
H_I^i(M)\otimes_RB\cong H_{IB}^i(M\otimes_RB).
\]
具体地，对域 \(k\)、\(R=k[x,y]\) 和 \(I=(x,y)\)，有
\[
H_I^0(R)=H_I^1(R)=0,\qquad
H_I^2(R)=\bigoplus_{a,b\ge1}k x^{-a}y^{-b},
\]
且其余次数为零。这里单项式表示它在 \(R_{xy}/(R_x+R_y)\) 中的商类，所列直和给出 \(k\) 向量空间的基，\(R\) 作用由这个余核继承。
\end{theorem}
\begin{proof}
双复形要把两种计算放在同一个总复形中：先计算内射消解方向，得到 \(M\) 的 Čech 复形；先计算 Čech 方向，则只剩定义中的挠子复形。

选内射消解 \(M\to E^\bullet\)，令
\[
\mathcal B^{p,q}=C^p(f_1,\ldots,f_r;R)\otimes_R E^q,
\qquad 0\le p\le r,\quad q\ge0.
\]
总次数为 \(p+q\)，在 \((p,q)\) 项上总微分为 \(\mathrm{d}_C+(-1)^p\mathrm{d}_E\)。先沿 \(q\) 方向取上同调：每个 Čech 项平坦，所以只剩 \(q=0\) 行的 \(C^\bullet(f_1,\ldots,f_r;M)\)。先沿 \(p\) 方向取上同调：上一个引理使每行只剩 \(p=0\) 列的 \(\Gamma_I(E^q)\)，从而得到定义中的复形。

两条边缘增广分别给出复形单射
\[
C^\bullet(f;M)\longrightarrow\operatorname{Tot}(\mathcal B),\qquad
\Gamma_I(E^\bullet)\longrightarrow\operatorname{Tot}(\mathcal B).
\]
对第一条取商后，每列底项换成 \(\mathcal B^{p,0}/(C^p(f;R)\otimes_RM)\)，其余项不变，各列正合。对第二条取商后，每行左端换成 \(\mathcal B^{0,q}/\Gamma_I(E^q)\)，其余项不变，各行正合。两个商总复形都正合，这可以直接作有限消元：列正合时，对次数 \(n\) 的圈选最小的非零横指标 \(p\)，则该分量 \(z_{p,q}\) 满足 \(\mathrm d_Ez_{p,q}=0\)。若 \(q=0\)，列底端正合已使它为零；否则选 \(w\) 使 \((-1)^p\mathrm d_Ew=z_{p,q}\)，减去总边界 \(\mathrm dw\) 便消去该分量，新增修正只在 \((p+1,q-1)\)。横指标最多到 \(n\)，所以有限步结束。行正合时交换两个指标作相同消元。由这两个正合商复形，边缘增广都诱导上同调同构，以上两种计算因而给出同一个总上同调。所有映射来自局部化及消解比较，所以同构自然。不同生成组的 Čech 复形可以有不同项数和微分；这里的生成组无关性指它们的上同调自然同构。

把 Čech 复形局部化于任一 \(f_j\)，其中的两项因子变成恒等映射 \(R_{f_j}\to R_{f_j}\)，可缩，故所有上同调局部化后为零。对于上同调中的一个元素，每个 \(f_j\) 的某个幂都将它零化。有限组生成元的任意足够高次乘积必含这样的幂，故 \(I\) 的某幂零化该元素。

对基变换公式，任意环同态 \(R\to B\) 都给出与局部化微分相容的逐项同构
\[
M_{f_{j_1}\cdots f_{j_p}}\otimes_RB
\cong(M\otimes_RB)_{f_{j_1}\cdots f_{j_p}}.
\]
因此 \(C^\bullet(f;M)\otimes_RB\cong C^\bullet(f;M\otimes_RB)\)，右边在 \(B\) 上计算。平坦性在下一步使用：张量保持微分的核、像及其商，所以
\(H^i(C^\bullet)\otimes_RB\cong H^i(C^\bullet\otimes_RB)\)。在两环上应用已证的 Čech 比较，就得到所述公式。

最后计算二元例子。将 \(R_x,R_y\) 看作 \(R_{xy}=k[x,x^{-1},y,y^{-1}]\) 的子环。零次核为零；在 \(R_x\) 中 \(y\) 的指数非负，在 \(R_y\) 中 \(x\) 的指数非负，故 \(R_x\cap R_y=R\)。一次复形的核因而恰为 \(R\) 的对角像，\(H_I^1(R)=0\)。二次余核 \(R_{xy}/(R_x+R_y)\) 商去了至少有一个指数非负的全部 Laurent 单项式；剩下的 \(x^{-a}y^{-b}\)、\(a,b\ge1\)，线性独立并张成该余核，所以构成所列基。复形长度为二，其余次数消失。
\end{proof}

例如 \(R=k[x]\)、\(I=(x)\) 时，\(H_I^0(R)=0\)，而
\[
H_I^1(R)=k[x,x^{-1}]/k[x]
=\bigoplus_{j\ge1}k x^{-j}.
\]
乘 \(x\) 把 \(x^{-j}\) 送到 \(x^{-j+1}\)，其中 \(x^{-1}\) 被送到零。每个元素被某个 \(x\) 的幂零化，但整个模不被统一的幂零化，也不有限生成。这是局部上同调与有限生成模上的 Ext 之间的一处显著差别。

\ProseParagraphBreak
这个模还具体记录了节首所见的满射失败。对短正合列
\[
0\longrightarrow R\xrightarrow{x}R\longrightarrow k\longrightarrow0,
\qquad R=k[x],
\]
关于 \((x)\) 的局部上同调长正合列含有
\[
0\longrightarrow k\xrightarrow{\delta}H^1_{(x)}(R)
\xrightarrow{x}H^1_{(x)}(R)\longrightarrow0.
\]
连接同态把 \(1\in k\) 送到 \(x^{-1}+R\)：先在中间的 \(R\) 中提升为 \(1\)，它的 Čech 微分是 \(1\in R_x\)，再沿前一项的乘 \(x\) 映射提升为 \(x^{-1}\)。因而 \(\ker(x:H^1_{(x)}(R)\to H^1_{(x)}(R))=k x^{-1}\)，恰好补上原来 \(\Gamma_{(x)}(R)=0\) 无法满射到 \(\Gamma_{(x)}(k)=k\) 的部分。高次局部上同调在这里把左正合函子的失效变成了可计算的正合列。

\subsection{局部上同调的消失与深度}
\label{b3:subsec:2201:03}

深度要求连续选取非零因子，局部上同调则测量支集在闭点的障碍。前几个上同调消失，恰好允许逐次选取相应长度的正则序列。

\begin{theorem}\label{b3:thm:depth-local-cohomology}
设 \(R\) 为交换 Noether 局部环，极大理想为 \(\mathfrak m\)，\(M\ne0\) 为有限生成的 \(R\) 模。则
\[
\operatorname{depth}_R M
=\min\{i\ge0:H_{\mathfrak m}^i(M)\ne0\}.
\]
\end{theorem}
\begin{proof}
商去一个正则元素使深度降一。长正合列将利用低次局部上同调上的乘法单射，结合逐元素的幂零化，使低次消失；再把商模的首次非零项送入原模高一次的上同调。

对 \(t=\operatorname{depth}_R M\) 归纳。若 \(t=0\)，\cref{b3:cor:depth-zero-socle}给出非零的 \(\operatorname{Soc}(M)\)，其每个元素被 \(\mathfrak m\) 零化，故 \(H_{\mathfrak m}^0(M)\ne0\)。

若 \(t>0\)，取 \(M\) 正则元素 \(x\in\mathfrak m\)。商 \(\bar M=M/xM\) 非零，且由\cref{b3:cor:regular-quotient-depth}有深度 \(t-1\)。原短正合列的局部上同调长正合列中含有
\[
H_{\mathfrak m}^{i-1}(\bar M)\longrightarrow
H_{\mathfrak m}^i(M)\xrightarrow{x}H_{\mathfrak m}^i(M).
\]
当 \(i<t\) 时左端为零，所以乘 \(x\) 在中间模上单射。该模为 \(\mathfrak m\)-挠子模；一个既被某个 \(x\) 的幂零化、又处在 \(x\) 单射的模中的元素只能为零。因此这些次数全部消失。取 \(i=t\)，长正合列给出非零模
\(H_{\mathfrak m}^{t-1}(\bar M)\) 到 \(H_{\mathfrak m}^t(M)\) 的单射，因为前一项 \(H_{\mathfrak m}^{t-1}(M)\) 已为零。于是第 \(t\) 次非零。
\end{proof}

\begin{theorem}\label{b3:thm:local-cohomology-dimension-vanishing}
设 \(R\) 为交换 Noether 局部环，极大理想为 \(\mathfrak m\)，\(M\ne0\) 为有限生成模，\(d=\dim_RM\)。则
\[
H_{\mathfrak m}^i(M)=0\quad(i>d).
\]
因此，\(M\) 为 Cohen--Macaulay 模当且仅当其局部上同调仅在次数 \(d\) 非零。
\end{theorem}
\begin{proof}
维数 \(d\) 允许选取 \(d\) 个参数，把支集切到闭点；这些参数给出的长度 \(d\) 的 Čech 复形就足以计算极大理想支集上的上同调，因此更高次数没有复形项。

取 \(M\) 的参数系 \(x_1,\ldots,x_d\)，其存在性由\cref{b3:thm:system-of-parameters}应用于 \(A=R/\operatorname{Ann}_RM\) 保证。在 \(A\) 中，所生成理想的根就是极大理想。Čech 复形的项只涉及 \(M\) 的局部化，所以不论把 \(M\) 看作 \(R\) 模还是 \(A\) 模，其复形相同。先在 \(A\) 上应用理想根无关性，再使用 Čech 比较定理，便可用 \(x_1,\ldots,x_d\) 的长度 \(d\) 复形计算 \(H_{\mathfrak m}^i(M)\)。更高次数没有项，因而消失。

若 \(M\) 为 CM 模，深度为 \(d\)，前一定理使低于 \(d\) 的次数消失并使第 \(d\) 次非零；反向同样由首次非零次数得到深度 \(d\)。
\end{proof}

对域 \(k\)，以 \(R=k[\![x,y]\!]/(x^2,xy)\) 为例，\((y)\) 的根为极大理想 \(\mathfrak m=(x,y)\)。每个元素唯一写成 \(a(y)+bx\)，其中 \(a(y)\in k[\![y]\!]\)、\(b\in k\)。由 \(yx=0\)，映射 \(R\to R_y=k(\!(y)\!)\) 的核是 \(kx\)，像是 \(k[\![y]\!]\)。因此
\[
H_{\mathfrak m}^0(R)=kx,\qquad
H_{\mathfrak m}^1(R)=k(\!(y)\!)/k[\![y]\!].
\]
环维数为一，深度却为零。被嵌入的零维部分 \(kx\) 在零次局部上同调中保留下来，一维主分量则在一次留下负幂的类，两部分信息出现在不同次数。CM 条件所要求的是全部局部上同调集中在模维数这一次数；下一节的对偶将说明，这种集中性为什么能用一个有限生成模表达。

% !TeX root = ../../Book3.tex
% 纲要标题：### 22.2 Matlis 对偶、典范模与局部对偶
% 纲要位置：计划纲要.md，第 2675 行。

\section{Matlis 对偶、典范模与局部对偶}
\label{b3:sec:2202}

局部上同调通常不是有限生成模；以环本身为值的 Hom 对偶甚至可能检测不到非零元素。内射包络与 Matlis 对偶提供合适的取值对象。若一个完备局部环来自完备正则局部环的满射表示，就能把它的局部上同调转为环境环上的 Ext；Cohen–Macaulay 性使对偶信息集中在一个次数，进而用典范模写出环自身上的对偶公式。

\ProseParagraphBreak
这些计算从满射 \(S\twoheadrightarrow R\) 出发，其中 \(S\) 是交换完备 Noether 正则局部环。幂级数商环天然带有这样的环境环；一般交换完备 Noether 局部环的相应满射来自\appref{b3:app:D}的 Cohen 结构定理。

% 纲要内容（仅作写作计划，尚非教材正文）：
%
% * 从离散赋值环的环值对偶失效，说明内射包络作为取值对象的作用
% * 内射包络、有限长度对偶与完备 Noether 局部环上的 Matlis 对偶；非完备双对偶反例及 Artin 不等于有限长度
% * 极小内射消解的基座计算；由 Koszul 对偶与内射模的极大理想幂挠部分计算正则局部环的最高局部上同调
% * 给定满射 $S\twoheadrightarrow R$，$S$ 为 $n$ 维完备 Noether 正则局部环、$d=\dim R$，构造典范模 $\omega_R=\operatorname{Ext}_S^{n-d}(R,S)$；此存在性不要求 $R$ 为 CM 环
% * 若 $R$ 为 CM 环，环境环上的 Ext 集中于 $n-d$ 次；换环引理给出单个典范模的局部对偶公式
% * 对有限生成的 $R$ 模比较局部上同调与 $\operatorname{Ext}_R^{d-i}(-,\omega_R)$ 的消失区间；非 CM 环以多个非零 Ext 次数说明单模公式的限制
% * 一般完备局部环上这类满射表示的存在性由附录 D 的 Cohen 结构定理说明
%
% ---
%

\subsection{环值对偶的不足}
\label{b3:subsec:2202:01}

设 \(k\) 为域，先看离散赋值环 \(A=k[\![t]\!]\)。由两项 Čech 复形得到
\[
H^1_{(t)}(A)=k(\!(t)\!)/k[\![t]\!].
\]
右端含有任意高的负幂，不能由有限个元素生成。普通的环值对偶在此为零。事实上，令 \(K=k(\!(t)\!)\)，则
\[
 \operatorname{Hom}_A(K/A,A)=0.
\]
这是因为每个 \(K/A\) 中的元素都被 \(t\) 的某个幂零化，而 \(A\) 没有非零的 \(t\)-挠元。因此以 \(A\) 为值的对偶无法检测这个非零模，需要改变取值对象。

\ProseParagraphBreak
这个对偶的取值对象将是剩余域的内射包络。它既能接受子模上同态的延拓，又能检测任意非零元素；在完备 Noether 局部环上，它把有限生成模与 Artin 模互相转换。因此，原来的局部上同调模虽然通常不有限生成，它的对偶却有可能由有限组生成元及关系描述。

\ProseParagraphBreak
给定满射 \(S\twoheadrightarrow R\)，其中 \(S\) 是 \(n\) 维交换完备 Noether 正则局部环，\(R\) 是 \(d\) 维商环，与最高次局部上同调对应的有限生成模是
\[
\operatorname{Ext}_S^{n-d}(R,S).
\]
它能够用 \(R\) 的有限长度、各项有限秩自由的 \(S\) 消解计算；下面的\cref{b3:def:canonical-module-regular-quotient}将从环自身的最高次局部上同调定义典范模，\cref{b3:prop:canonical-module-quotient}再识别它与这个环境 Ext 模。余维 \(n-d\) 出现在这里，是因为正则局部环的最高局部上同调位于次数 \(n\)，而 \(R\) 的维数为 \(d\)。若 \(R\) 是 Cohen–Macaulay 环，其余局部上同调消失，整个对偶公式便集中到这一个模上。

\subsection{离散赋值环的内射对偶}
\label{b3:subsec:2202:02}

先把上述离散赋值环的例子算得更具体。令 \(A=k[\![t]\!]\)、\(K=k(\!(t)\!)\)、\(E=K/A\)。第二卷第7.2节的 Baer 判据说明 \(E\) 是内射 \(A\) 模：任意非零理想为 \(t^rA\)，若 \(f(t^r)=a+A\)，则取 \(g(1)=t^{-r}a+A\)，就得到延拓 \(g:A\rightarrow E\)。零理想的映射用零映射延拓。

\ProseParagraphBreak
\(E\) 中每个非零元素都能乘以某个 \(t\) 的幂，得到非零的 \(t^{-1}\) 项。因此子模 \(kt^{-1}\cong k\) 与 \(E\) 的任意非零子模相交非零。这说明内射性没有加入与剩余域毫无关系的直和因子；这正是下文“内射包络”中包络条件的含义。

\ProseParagraphBreak
记 \(D_A(N)=\operatorname{Hom}_A(N,E)\)。这里有
\[
D_A(A)=E,
\qquad
D_A(A/t^rA)=(0:_Et^r)=t^{-r}A/A\cong A/t^rA.
\]
最后一个同构把 \(a+t^rA\) 送到 \(at^{-r}+A\)。这使有限长度模的对偶仍是有限长度模，而自由模的对偶是非有限生成的 \(E\)。离散赋值环的计算对应于一般局部对偶公式
\[
H^i_{\mathfrak m}(M)
\cong D_R\bigl(\operatorname{Ext}_R^{d-i}(M,\omega_R)\bigr),
\]
其中 \(R=S/I\) 是来自完备 Noether 正则局部环 \(S\) 的 \(d\) 维 Cohen–Macaulay 商环，\(\mathfrak m\) 是其极大理想，\(M\) 是有限生成的 \(R\) 模；\(\omega_R\) 是将在\cref{b3:def:canonical-module-regular-quotient}定义的典范模。完整陈述和证明见\cref{b3:thm:canonical-module-local-duality}。

\subsection{局部对偶的适用条件}
\label{b3:subsec:2202:03}

局部对偶中的三种条件各有作用。Noether 性保证有限生成模有有限展示，并使内射包络的元素能用极大理想的幂控制。完备性使取两次对偶后的环仍是原环；若环不完备，首先出现的是它的完备化。Cohen–Macaulay 性则把环自身的局部上同调集中在唯一的次数 \(d\)，使这些对偶公式都可以用一个模 \(\omega_R\) 表示。

\ProseParagraphBreak
这些条件不能由“局部环”三个字替代。一般非 Cohen–Macaulay 环可能同时有几个非零局部上同调模，不能把上述公式中的 \(\omega_R\) 原样用于所有次数。即使某个环具有典范模，也不表示它的全部对偶信息已经集中在该模上。\subsecref{b3:subsec:2202:06}将给出一个含有嵌入点的幂级数商环，直接检验这种失败。

\ProseParagraphBreak
若只给出一个交换完备 Noether 局部环 \(R\)，Cohen 结构定理提供来自完备正则局部环的满射表示；相应的系数域与系数环构造见\appref{b3:app:D}。

\subsection{内射包络与有限生成模、Artin 模的 Matlis 对偶}
\label{b3:subsec:2202:04}

先说明如何在所有可能的内射扩张中选择一个尽量小的对象。这里使用第二卷第7.2节已证明的两个结果：每个模都可嵌入内射模；Baer 判据用理想上的同态延拓检验内射性。内射包络及以下对偶定理在本节证明。松村英之《Commutative Ring Theory》的对应论述见\cite[\S 18, Theorems 18.4--18.7, pp. 145--150]{Matsumura1989CommutativeRingTheory}。

\begin{definition}\label{b3:def:injective-hull}
设 \(R\) 为交换环，\(N\subseteq M\) 为 \(R\) 模。若 \(M\) 的每个非零子模都与 \(N\) 相交非零，则此包含为\term[benzhi-kuozhang]{本质扩张}[essential extension]。若 \(M\) 还为内射模，就称 \(N\subseteq M\) 为 \(N\) 的\term[neishe-baoluo]{内射包络}[injective hull]，记一个这样的模为 \(E_R(N)\)。
\end{definition}
\symidx{\(E_R(N)\)：模 \(N\) 的内射包络}

本质性有一个常用的检验形式：若映射 \(f:M\rightarrow L\) 在 \(N\) 上单射，则它在整个 \(M\) 上单射，因为非零的 \(\ker f\) 必与 \(N\) 相交。

\ProseParagraphBreak
本质扩张要求每个非零子模都遇到原来的 \(N\)，所以不能额外加入一个与 \(N\) 交零的非零直和因子。前面的 \(kt^{-1}\subset k(\!(t)\!)/k[\![t]\!]\) 就体现了这一点：任意非零子模都有元素经乘法落到这份剩余域中。下面的存在性证明先在一个内射模内取极大本质扩张，再证明它本身内射；唯一性同样由本质性排除补因子。

\begin{proposition}\label{b3:prop:injective-hull-existence}
设 \(R\) 为交换环，\(N\) 为 \(R\) 模。\(N\) 有内射包络，且任意两个内射包络之间存在限制在 \(N\) 上为恒等的同构。这里不声称这个同构唯一。
\end{proposition}
\begin{proof}
把 \(N\) 嵌入内射模 \(J\)。在 \(J\) 中所有本质包含 \(N\) 的子模中按包含排序；一条链的并仍本质包含 \(N\)，因为并中的非零元素已落在链的某个成员中。Zorn 引理给出极大者 \(E\)。

\(E\) 没有真本质扩张。事实上，若 \(E\subseteq X\) 本质，则本质性的传递性使 \(N\subseteq X\) 本质；内射性把 \(E\hookrightarrow J\) 延拓为 \(X\rightarrow J\)。此映射的核与 \(E\) 交零，故为零。其像在 \(J\) 中仍本质包含 \(N\)，极大性使像等于 \(E\)。映射又在 \(E\) 上为恒等，故 \(X=E\)。

再把 \(E\) 嵌入内射模 \(J'\)，选取与 \(E\) 交零的极大子模 \(L\subseteq J'\)。商模 \(J'/L\) 本质包含 \(E\)：若一个非零子模与 \(E\) 的像交零，其原像就是真包含 \(L\) 且与 \(E\) 交零的子模，矛盾。由上一段，\(J'/L=E\) 的像，故 \(J'=E\oplus L\)。内射模的直和因子内射，得存在性。

若 \(E,E'\) 都是包络，由 \(E'\) 的内射性延拓 \(N\) 上的恒等为 \(u:E\rightarrow E'\)。本质性使 \(u\) 单射，而内射子模 \(u(E)\) 在 \(E'\) 中分裂。补因子与 \(N\) 交零，本质性迫使补因子为零，所以 \(u\) 是同构。
\end{proof}

以下固定交换 Noether 局部环 \((R,\mathfrak m,k)\) 及包络 \(k\subseteq E=E_R(k)\)。

\ProseParagraphBreak
其基座只有指定的一份剩余域，每个元素又被极大理想的某个幂零化。这两点控制 \(E\) 中的有限层；第三项元素检测性质则保证，以 \(E\) 为值的对偶不会把非零模全部送到零。

\begin{proposition}\label{b3:prop:injective-hull-local-properties}
设 \((R,\mathfrak m,k)\) 为交换 Noether 局部环，\(E=E_R(k)\)。则
\[
\operatorname{Soc}_R E=(0:_E\mathfrak m)=k,
\qquad
E=\bigcup_{r\ge1}(0:_E\mathfrak m^r).
\]
此外，对任意 \(R\) 模 \(M\) 及 \(0\ne x\in M\)，存在 \(f:M\rightarrow E\) 使 \(f(x)\ne0\)。
\end{proposition}
\begin{proof}
\((0:_E\mathfrak m)\) 是 \(k\) 向量空间；它的每条一维子空间都须与指定的 \(k\) 相交，所以它恰为该 \(k\)。

固定 \(e\in E\)。若 \(e=0\)，它已被 \(\mathfrak m\) 零化；以下设 \(e\ne0\)。若 \(\mathfrak p\in\operatorname{Ass}_R(Re)\)，则 \(Re\) 有一个同构于 \(R/\mathfrak p\) 的子模。该子模与 \(k\) 相交非零；而 \(R/\mathfrak p\) 的每个非零元素的零化子都是 \(\mathfrak p\)，相交元素的零化子又是 \(\mathfrak m\)，故 \(\mathfrak p=\mathfrak m\)。对由 \(e\) 生成的循环模 \(Re\) 应用\cref{b3:prop:minimal-primes-associated}，得到 \(\sqrt{\operatorname{Ann}_R(e)}=\mathfrak m\)。\(\mathfrak m\) 有限生成，因而某个 \(\mathfrak m^r\) 包含于该零化子，得并集公式。

对非零 \(x\)，\(\operatorname{Ann}_R(x)\) 是真理想。复合
\(Rx\cong R/\operatorname{Ann}_R(x)\twoheadrightarrow k\hookrightarrow E\)
把 \(x\) 送到非零元素，再用内射性将它延拓到 \(M\)。
\end{proof}

\begin{definition}\label{b3:def:matlis-dual}
设 \((R,\mathfrak m,k)\) 为交换 Noether 局部环，固定 \(E=E_R(k)\)。对 \(R\) 模 \(M\)，模
\[
D_R(M)\coloneqq\operatorname{Hom}_R(M,E)
\]
为它的\term[Matlis-duiou]{Matlis 对偶}[Matlis dual]。同态 \(u:M\rightarrow N\) 对应预复合映射 \(D_R(u):D_R(N)\rightarrow D_R(M)\)。
\end{definition}
\symidx{\(D_R(M)\)：模 \(M\) 的 Matlis 对偶}

\(D_R\) 是反变正合函子，因为 \(E\) 内射。上一命题还说明 \(D_R(M)=0\) 蕴含 \(M=0\)。求值映射
\[
\delta_M:M\longrightarrow D_RD_R(M),
\qquad \delta_M(x)(f)=f(x)
\]
对每个模都单射。不过，满射性需要限制模的类别。

\begin{lemma}\label{b3:lem:matlis-finite-length}
设 \((R,\mathfrak m,k)\) 为交换 Noether 局部环，\(D_R(-)=\operatorname{Hom}_R(-,E_R(k))\)。若 \(M\) 为有限长度 \(R\) 模，\(\delta_M:m\mapsto(\varphi\mapsto\varphi(m))\) 为求值映射，则
\[
\ell_R D_R(M)=\ell_R M,
\qquad \delta_M:M\xrightarrow{\sim}D_RD_R(M).
\]
并且乘法给出环同构
\[
\widehat R\xrightarrow{\sim}\operatorname{End}_R\bigl(E_R(k)\bigr),
\qquad \widehat R=\varprojlim_r R/\mathfrak m^r.
\]
具体地，若 \(k\) 为域、\(R_0=k[t]_{(t)}\)、\(E_0=k(t)/R_0\)、\(D_{R_0}=\operatorname{Hom}_{R_0}(-,E_0)\)，则 \(E_0\) 为 \(k\) 的内射包络，且
\[
D_{R_0}D_{R_0}(R_0)\cong k[\![t]\!].
\]
自然求值映射对应完备化映射 \(R_0\to k[\![t]\!]\)，它不是满射。
\end{lemma}
\begin{proof}
先由 \(D_R(k)=\operatorname{Soc}E\cong k\) 及正合性，对组成列长度归纳得到长度等式。求值映射已知单射，且两端长度相等，故为同构。

令 \(E_r=(0:_E\mathfrak m^r)=D_R(R/\mathfrak m^r)\)。前述有限长度对偶给出
\(D_R(E_r)\cong R/\mathfrak m^r\)，且元素 \(a+\mathfrak m^r\) 对应 \(e\mapsto ae\)。由 \(E=\bigcup_rE_r\)，一个 \(E\) 上的同态等价于各 \(E_r\) 上相容的同态族。因此
\[
\operatorname{End}_R(E)
=\varprojlim_r\operatorname{Hom}_R(E_r,E)
\cong\varprojlim_rR/\mathfrak m^r.
\]
这些同构把复合变为乘法，故是环同构。

对具体例子，\(R_0=k[t]_{(t)}\) 是 DVR，每个非零理想主生成，而 \(E_0=k(t)/R_0\) 可除。对主理想上同态，除以其生成元就能延拓到 \(R_0\)，Baer 判据因而给出内射性。其基座是 \(kt^{-1}\)，任意非零元素取适当的 \(t\) 幂后都成为这个基座中的非零元素，所以 \(k\subseteq E_0\) 本质，\(E_0\) 是内射包络。\(D_{R_0}(R_0)=E_0\)，故前面已证的自同态计算给出
\[
D_{R_0}D_{R_0}(R_0)=\operatorname{End}_{R_0}(E_0)
\cong\widehat R_0=k[\![t]\!].
\]
求值作用在 \(E_0\) 上就是乘法，因而自然映射正是完备化映射。它不满：级数 \(f(t)=\sum_{j\ge1}t^{j!}\) 不在 \(R_0\) 中。若 \(f=p/q\)，其中 \(p,q\in k[t]\)、\(q(0)\ne0\)，取 \(j\) 足够大，使 \(j!>\deg p\) 且 \(j!-(j-1)!>\deg q\)。在 \(qf=p\) 中比较 \(t^{j!}\) 的系数，左边只有 \(q(0)\ne0\)，右边为零，矛盾。
\end{proof}

完备性在这里有了具体作用：只有 \(R=\widehat R\) 时，\(D_RD_R(R)=D_R(E)\) 才由同一个环 \(R\) 表示。以下的\term[Matlis-duiou-dingli]{Matlis 对偶定理}[Matlis duality theorem]将这一等式推广到两个完整的模类别。

\begin{theorem}[Matlis 对偶]\label{b3:thm:matlis-duality}
设 \((R,\mathfrak m,k)\) 为交换完备 Noether 局部环，\(E=E_R(k)\)，\(D_R(-)=\operatorname{Hom}_R(-,E)\)。则 \(E\) 是 Artin 模；\(D_R\) 把有限生成的 \(R\) 模变成 \(R\) 上的 Artin 模，把 \(R\) 上的 Artin 模变成有限生成的 \(R\) 模。对这两类中的任意模 \(M\)，求值映射 \(\delta_M:M\to D_RD_R(M)\)、\(m\mapsto(\varphi\mapsto\varphi(m))\) 是同构。因此 \(D_R\) 在有限生成模的范畴与 Artin 模的范畴之间给出反变等价。
\end{theorem}
\begin{proof}
证明先将 \(E\) 的子模降链转成 \(R\) 的理想升链，得到 \(E\) 的 Artin 性。有限生成模的对偶随后嵌入有限份 \(E\)，Artin 模本身也将嵌入有限份 \(E\)，但后一项须借助降链条件和元素检测。有限展示给出有限生成一侧的双对偶同构，再用求值映射的相容关系返回 Artin 一侧。

先证 \(E\) 满足降链条件。对子模 \(L\subseteq E\) 记
\[
L^\perp=\{\,a\in R\mid aL=0\,\},
\qquad
J^\perp=(0:_EJ)\quad(J\subseteq R\;\text{为理想}).
\]
这里 \(L^\perp\) 是 \(R\) 中的理想，而 \(J^\perp\) 是 \(E\) 中的子模，两次取垂直才回到原来的对象类型。
显然 \(L\subseteq L^{\perp\perp}\)。若 \(e\notin L\)，元素检测性质在 \(E/L\) 上给出同态 \(f:E/L\rightarrow E\)，使 \(f(e+L)\ne0\)。其与商映射的复合是 \(E\) 的自同态，前一引理和完备性使它为乘以某个 \(a\in R\)。此时 \(aL=0\) 而 \(ae\ne0\)，所以 \(e\notin L^{\perp\perp}\)。故 \(L=L^{\perp\perp}\)。同样，对 \(a\notin J\)，在 \(R/J\) 上检测 \(a+J\) 的同态给出 \(e\in J^\perp\) 且 \(ae\ne0\)，故 \(J=J^{\perp\perp}\)。在当前完备 Noether 局部环的假设下，这两个取垂直的操作互为逆，且反转包含关系：
\[
\begin{tikzcd}[column sep=5em]
\{L\subseteq E:L\;\text{为子模}\}
\arrow[r, shift left=1.2ex, "L\mapsto L^\perp"] &
\{J\subseteq R:J\;\text{为理想}\}
\arrow[l, shift left=1.2ex, "J\mapsto J^\perp"]
\end{tikzcd}
\]
\[
L^{\perp\perp}=L,\qquad J^{\perp\perp}=J,\qquad
L_1\subseteq L_2\ \Longleftrightarrow\ L_2^\perp\subseteq L_1^\perp.
\]
因此 \(E\) 的子模降链变成 \(R\) 的理想升链；后者稳定，取垂直后前者也稳定，故 \(E\) 为 Artin 模。

若 \(M\) 有限生成，由满射 \(R^r\twoheadrightarrow M\) 得单射 \(D_R(M)\hookrightarrow E^r\)，故 \(D_R(M)\) 为 Artin 模。\(R\) 的 Noether 性又给出有限展示 \(R^s\rightarrow R^r\rightarrow M\rightarrow0\)。\(D_RD_R\) 是正合协变函子，\(\delta_R\) 由前一引理是同构，所以将该展示取两次对偶并比较余核，得到 \(\delta_M\) 同构。

若 \(A\) 为 Artin 模，考察全部映射 \(A\rightarrow E^r\) 的核，其中 \(r\) 可为任意非负整数。降链条件保证这些核中有极小者 \(K\)。若 \(0\ne x\in K\)，可再加一个分量 \(A\rightarrow E\) 检测 \(x\)，把核严格缩小，矛盾。因此某个映射给出单射 \(A\hookrightarrow E^r\)。取对偶得到满射 \(R^r\twoheadrightarrow D_R(A)\)，故 \(D_R(A)\) 有限生成。

最后，直接计算求值映射可得
\[
D_R(\delta_A)\circ\delta_{D_R(A)}
=1_{D_R(A)}.
\]
因为 \(D_R(A)\) 有限生成，\(\delta_{D_R(A)}\) 已是同构，故 \(D_R(\delta_A)\) 同构。\(D_R\) 正合且检测非零模，遂使 \(\delta_A\) 的核和余核均为零。求值映射的自然性说明这些同构与一切模同态相容，给出所述反变等价。
\end{proof}

Artin 性与有限长度仍须区分。设 \(k\) 为域、\(A=k[\![t]\!]\)，则 \(E=k(\!(t)\!)/A\) 由上述定理为 Artin 模，满足子模降链条件；但它包含严格上升的子模链 \(t^{-1}A/A\subsetneq t^{-2}A/A\subsetneq\cdots\)，所以不是 Noether 模，更不是有限长度模。有限长度要求升链、降链同时稳定，Matlis 对偶的 Artin 一侧允许这里无界的分母。

\subsection{典范模及 Cohen–Macaulay 情形}
\label{b3:subsec:2202:05}

正则局部环的最高次局部上同调给出典范模构造中的对偶对象。计算时，内射消解中被极大理想的幂零化的部分由基座决定；极小内射消解又把各项的基座直接转为剩余域上的 Ext。

\begin{lemma}\label{b3:lem:injective-maximal-torsion}
设 \((A,\mathfrak n,k)\) 为交换 Noether 局部环，\(E=E_A(k)\)，\(J\) 为内射 \(A\) 模。若 \(\{u_\lambda\}_{\lambda\in\Lambda}\) 为 \(\operatorname{Soc}_A J\) 的一个 \(k\) 基，则
\[
\Gamma_{\mathfrak n}(J)\cong E^{(\Lambda)}.
\]
其中 \(E^{(\Lambda)}\) 为直和，指标数等于该基座的维数。
\end{lemma}
\begin{proof}
先注意，在 Noether 环上内射模的任意直和仍内射。由 Baer 判据，只需把理想 \(L\) 上的映射延拓到 \(A\)。因为 \(L\) 有限生成，其像落在有限个直和项中；逐分量延拓后再合成即可。

把每个映射 \(k\rightarrow J\)、\(1\mapsto u_\lambda\) 延拓到 \(E\rightarrow J\)，并合成为 \(v:E^{(\Lambda)}\rightarrow J\)。直和的每个元素由有限个分量组成，且由\cref{b3:prop:injective-hull-local-properties}被某个 \(\mathfrak n\) 的幂零化。因此它的任何非零子模都含有被 \(\mathfrak n\) 零化的非零元素：取零化幂数最小的元素，再乘 \(\mathfrak n\) 的适当幂。故 \(k^{(\Lambda)}\) 在该直和中本质。\(v\) 在这一个基座上单射，遂整体单射。

其像内射，故 \(J=v(E^{(\Lambda)})\oplus J'\)。\(J'\) 的基座为零，因为 \(u_\lambda\) 已张成 \(J\) 的全部基座。若 \(\Gamma_{\mathfrak n}(J')\ne0\)，同样取最小零化幂数就会找到非零基座元素，矛盾。因此极大理想幂挠部分恰为第一直和项。
\end{proof}

\begin{lemma}\label{b3:lem:minimal-injective-resolution-socle}
设 \((A,\mathfrak n,k)\) 为交换 Noether 局部环，\(M\) 为 \(A\) 模。令 \(C^0=M\)，对 \(q\ge0\) 递归取内射包络 \(C^q\subseteq J^q=E_A(C^q)\)、商模 \(C^{q+1}=J^q/C^q\)。以商映射与下一次包络的包含复合为微分，所得 \(M\to J^\bullet\) 是 \(M\) 的极小内射消解，且
\[
\operatorname{Hom}_A(k,J^\bullet)\;\text{的微分均为零},
\qquad
\operatorname{Soc}_A J^q\cong\operatorname{Ext}_A^q(k,M).
\]
若某个整数 \(n\ge0\) 满足 \(\operatorname{Ext}_A^{n+1}(A/L,M)=0\) 对所有理想 \(L\subseteq A\) 成立，则 \(C^n\) 内射，可在第 \(n\) 项结束此消解。
\end{lemma}
\begin{proof}
第十七章的极小自由消解张量剩余域后微分为零。极小内射消解相应地给出零映射
\[
\operatorname{Hom}_A(k,J^q)\xrightarrow{0}
\operatorname{Hom}_A(k,J^{q+1})\qquad(q\ge0).
\]
前者逐层选择极小生成组，后者逐层取本质的内射包络，使每个简单子模都留在当前微分的核中。

由构造，\(J^q\to J^{q+1}\) 的核是 \(C^q\)，在 \(q\ge1\) 时又等于前一微分的像，所以增广列正合。若 \(u\in\operatorname{Soc}_A J^q\) 非零，则 \(Au\cong k\) 是简单子模。本质性使 \(Au\cap C^q\ne0\)，故 \(Au\subseteq C^q\)。因此全部基座元素都在微分的核中，取 \(\operatorname{Hom}_A(k,-)\) 后微分为零。

内射消解计算 Ext 与自由消解计算 Ext 给出同一个模。具体地，对自由消解 \(P_\bullet\to k\)，取双复形 \(\operatorname{Hom}_A(P_p,J^q)\)，其中 \(p,q\ge0\)、总次数为 \(p+q\)。沿 \(q\) 方向，\(P_p\) 的投射性使 Hom 保持消解正合，只留下 \(q=0\) 行的 \(\operatorname{Hom}_A(P_\bullet,M)\)；沿 \(p\) 方向，\(J^q\) 的内射性使正次数消失，只留下 \(p=0\) 列的 \(\operatorname{Hom}_A(k,J^\bullet)\)。两个消解虽可无限延伸，第一象限中的每条固定总次数对角线仍只有有限项，逐次提升圈并消去正合行、列中的分量便得到两种计算之间的同构。后一个复形微分为零，故其第 \(q\) 次上同调就是 \(\operatorname{Hom}_A(k,J^q)=\operatorname{Soc}_A J^q\)。

最后，短正合列 \(0\to C^q\to J^q\to C^{q+1}\to0\) 逐次作维数平移，给出
\[
\operatorname{Ext}_A^1(A/L,C^n)
\cong\operatorname{Ext}_A^{n+1}(A/L,M)=0.
\]
Baer 判据使 \(C^n\) 内射。它既内射又在 \(J^n\) 中本质，便只能是 \(J^n\) 自身，于是 \(C^{n+1}=0\)。当 \(n=0\) 时，维数平移等式不需经过中间项，结论仍成立。
\end{proof}

\begin{proposition}\label{b3:prop:regular-local-top-cohomology}
设 \((S,\mathfrak n,k)\) 为 \(n\) 维交换 Noether 正则局部环。则
\[
H^q_{\mathfrak n}(S)=0\quad(q\ne n),
\qquad
H^n_{\mathfrak n}(S)\cong E_S(k).
\]
此处不要求 \(S\) 完备，最后的同构依赖于一个剩余域基座的识别。
\end{proposition}
\begin{proof}
先用正则参数的 Koszul 对偶计算 \(\operatorname{Ext}_S^q(k,S)\)，再由极小内射消解把这些 Ext 识别为各项基座。内射模的极大理想幂挠部分由基座决定，因而最终只在最高次数留下剩余域的内射包络。

取 \(\mathfrak n\) 的正则生成系 \(x_1,\ldots,x_n\)，其存在由\cref{b3:thm:regular-local-characterizations}保证。\cref{b3:thm:koszul-regular-resolution}给出 \(k\) 的 Koszul 自由消解。外积配对
\(\bigwedge^qS^n\times\bigwedge^{n-q}S^n\rightarrow\bigwedge^nS^n\cong S\)
把其对偶复形识别为反向的 Koszul 复形。具体地，令 \(\Phi_q(v)(u)\) 为 \(u\wedge v\) 的最高外积系数。对 \(u\in K_{q+1}\)、\(v\in K_{n-q}\)，等式 \(\mathrm{d}(u\wedge v)=0\) 给
\(\delta\Phi_q(v)=(-1)^q\Phi_{q+1}(\mathrm{d}v)\)。故以 \((-1)^{q(q-1)/2}\Phi_q\) 作第 \(q\) 项的识别，所有微分均相容。因此
\[
\operatorname{Ext}_S^q(k,S)=
\begin{cases}k,&q=n,\\0,&q\ne n.\end{cases}
\]
等价地，也可将 \(n\) 个两项复形 \(S\xrightarrow{x_i}S\) 逐个对偶，每对偶一个因子便把它的非零余核提高一个次数。

取\cref{b3:lem:minimal-injective-resolution-socle}构造的 \(S\) 的极小内射消解 \(J^\bullet\)。由该引理，\(\operatorname{Soc}_S J^q\) 的维数在 \(q=n\) 时为一，在其余次数为零。又由\cref{b3:thm:regular-local-global-dimension}，对所有理想 \(L\subseteq S\) 都有 \(\operatorname{Ext}_S^{n+1}(S/L,S)=0\)，所以同一引理保证消解在第 \(n\) 项结束。

\cref{b3:lem:injective-maximal-torsion}于是给出
\[
\Gamma_{\mathfrak n}(J^q)\cong
\begin{cases}E_S(k),&q=n,\\0,&q\ne n.\end{cases}
\]
这个复形只有一个非零项。按支集挠子函子的导出函子计算局部上同调，即得结论；这也包括 \(n=0\)、\(S=k\) 的情形。
\end{proof}

\begin{theorem}\label{b3:thm:regular-local-duality}
设 \((S,\mathfrak n,k)\) 为 \(n\) 维交换完备 Noether 正则局部环，\(E=E_S(k)\)，\(D_S(-)=\operatorname{Hom}_S(-,E)\)。选定同构 \(H^n_{\mathfrak n}(S)\cong E\) 后，对每个有限生成的 \(S\) 模 \(N\) 和每个整数 \(i\)，有关于 \(N\) 自然的同构
\[
H^i_{\mathfrak n}(N)
\cong D_S\bigl(\operatorname{Ext}_S^{n-i}(N,S)\bigr),
\qquad
D_S\bigl(H^i_{\mathfrak n}(N)\bigr)
\cong\operatorname{Ext}_S^{n-i}(N,S).
\]
负次局部上同调和负次 Ext 均约定为零。各 \(H^i_{\mathfrak n}(N)\) 均为 Artin 模。
\end{theorem}
\begin{proof}
要比较局部上同调与环境 Ext，可以用同一个双复形的两种计算。自由消解方向先计算，会返回 \(N\) 的 Čech 上同调；Čech 方向先计算，则由前一命题只留下最高次数的内射包络，再用有限秩自由性把所得复形识别为 Ext 计算的 Matlis 对偶。

令 \(C^\bullet\) 为正则生成系 \(x_1,\ldots,x_n\) 的 Čech 复形。\cref{b3:thm:regular-local-homological-criterion}保证 \(N\) 有长度至多 \(n\)、各项有限秩自由的消解 \(P_\bullet\rightarrow N\)。将 \(P_j\) 放在上链次数 \(-j\)，取双复形
\[
C^q\otimes_SP_j,\qquad 0\le q\le n,\quad 0\le j\le n,
\]
其总次数为 \(q-j\)。总微分采用
\(\mathrm{d}(c\otimes p)=\mathrm{d}_Cc\otimes p+(-1)^{\deg c}c\otimes \mathrm{d}_Pp\)，故平方为零。

每个 \(C^q\) 都平坦，所以沿 \(P\) 方向取同调只在 \(j=0\) 留下 \(C^q\otimes_SN\)。因此总上同调为 \(H^i_{\mathfrak n}(N)\)，这里使用\cref{b3:thm:cech-local-cohomology}。换一个方向，前一命题给 \(H^q(C^\bullet)=0\) 对 \(q\ne n\)，而 \(H^n(C^\bullet)=E\)。各 \(P_j\) 自由，故这次只留下 \(q=n\) 行的 \(E\otimes_SP_j\)。该行的微分为 \((-1)^n\mathrm{d}_P\)，在第 \(j\) 项乘以 \((-1)^{nj}\) 就将其识别为通常的链复形 \(E\otimes_SP_\bullet\)。两次比较用有限滤过完成：按一方向次数过滤总复形，把圈从最外侧的分量开始处理，利用正合性减去边界，有限次后只剩非零同调所在的行；边界的比较同样进行。两个方向均有界，故过程终止，并与复形映射相容。

最高行位于 Čech 次数 \(n\)，自由消解的链次数 \(j\) 将总次数降到 \(n-j\)，所以在总次数 \(i\) 必有 \(j=n-i\)。这给出
\[
H^i_{\mathfrak n}(N)
\cong H_{n-i}(E\otimes_SP_\bullet).
\]
有限自由性给出逐项同构
\[
E\otimes_SP_j
\cong\operatorname{Hom}_S\bigl(\operatorname{Hom}_S(P_j,S),E\bigr),
\quad e\otimes p\longmapsto[\lambda\mapsto\lambda(p)\cdot e].
\]
此同构与微分相容。\(D_S\) 正合，故右侧复形的同调为 \(D_S(\operatorname{Ext}_S^{n-i}(N,S))\)，得到第一个同构。Ext 模有限生成，Matlis 对偶定理使其对偶为 Artin 模；再取一次对偶并使用有限生成模的求值同构，得到第二个同构。
\end{proof}

现在可以准确地命名所需的有限生成模。

\begin{definition}\label{b3:def:canonical-module-regular-quotient}
设 \((R,\mathfrak m,k)\) 为 \(d\) 维交换完备 Noether 局部环，\(D_R=\operatorname{Hom}_R(-,E_R(k))\)。若有限生成的 \(R\) 模 \(\omega_R\) 满足
\[
D_R(\omega_R)\cong H^d_{\mathfrak m}(R),
\]
则称 \(\omega_R\) 为 \(R\) 的\term[dianfanmo]{典范模}[canonical module]。
\end{definition}
\symidx{\(\omega_R\)：局部环 \(R\) 的典范模}

Matlis 对偶定理说明，一旦典范模存在，对定义中的同构再取对偶便得
\[
\omega_R\cong D_R\bigl(H^d_{\mathfrak m}(R)\bigr).
\]
这从 \(R\) 自身的最高次局部上同调确定典范模的同构类；下面的环境 Ext 则给出它的一个具体实现。这里不指定某一个同构，所以“典范”不表示所有实现之间有唯一的同构，与内射包络的唯一性一样，它确定的是同构类。

\begin{proposition}\label{b3:prop:canonical-module-quotient}
设 \((S,\mathfrak n,k)\) 为 \(n\) 维交换完备 Noether 正则局部环，\(I\subsetneq S\) 为理想，\(R=S/I\)、\(\mathfrak m=\mathfrak n/I\)、\(d=\dim R\)。则
\[
\omega_R=\operatorname{Ext}_S^{n-d}(R,S)
\]
是 \(R\) 的典范模。若 \(R\) 还是 Cohen–Macaulay 环，令 \(c=n-d\)，则
\[
\operatorname{Ext}_S^q(R,S)=0\quad(q\ne c),
\qquad
\operatorname{depth}_R\omega_R=d.
\]
从不同的完备正则局部环到 \(R\) 的满射表示产生同构的典范模。
\end{proposition}
\begin{proof}
先比较两个环上的对偶对象。模
\[
E_R=\operatorname{Hom}_S(R,E_S(k))=(0:_{E_S(k)}I)
\]
作为 \(R\) 模内射，因为对每个 \(R\) 模 \(M\)，在 \(1\) 处求值给出自然同构
\[
\operatorname{Hom}_R(M,E_R)\cong\operatorname{Hom}_S(M,E_S(k)).
\]
右端是正合函子。\(E_R\) 中的 \(k\) 本质，因为它已经在 \(E_S(k)\) 中本质。因此 \(E_R\) 就是 \(E_R(k)\) 的一个实现，而 \(D_R(M)=D_S(M)\)。由同一组生成元的 Čech 复形，\(H^i_{\mathfrak m}(M)=H^i_{\mathfrak n}(M)\)。正则局部对偶应用于 \(M=R\)，便给出典范模定义要求的同构；\cref{b3:prop:finite-ext-localization,b3:lem:ext-basic-properties}分别说明 Ext 有限生成且被第一变量的零化子 \(I\) 零化，故是有限生成的 \(R\) 模。

现假设 \(R\) 为 CM 环。\cref{b3:thm:depth-local-cohomology}与\cref{b3:thm:local-cohomology-dimension-vanishing}使 \(H^i_{\mathfrak m}(R)\) 仅在 \(i=d\) 非零。正则局部对偶立即给出 Ext 的集中次数。

进一步，\(R\) 作为 \(S\) 模的深度也是 \(d\)，因为 \(\mathfrak n\) 中元素在 \(R\) 上的作用只由其在 \(\mathfrak m\) 中的像决定。由\cref{b3:thm:auslander-buchsbaum}，\(\operatorname{pd}_SR=c\)。取长度 \(c\)、各项有限秩自由的消解 \(F_\bullet\rightarrow R\)。对偶复形 \(\operatorname{Hom}_S(F_\bullet,S)\) 仅在最后一项有上同调 \(\omega_R\)；反向排列后，它就是 \(\omega_R\) 的长度 \(c\) 的自由消解。再对偶一次，逐项双对偶回到 \(F_\bullet\)，所以
\[
\operatorname{Ext}_S^q(\omega_R,S)=
\begin{cases}R,&q=c,\\0,&q\ne c.\end{cases}
\]
故 \(\operatorname{pd}_S\omega_R=c\)，再次应用 Auslander–Buchsbaum 公式得深度 \(n-c=d\)。\(S\) 与 \(R\) 上的深度相同，得到所述结论。不同满射表示所构造的典范模具有相同的同构类型，因为它们都同构于 \(D_R(H^d_{\mathfrak m}(R))\)。
\end{proof}

存在性一项本身没有要求 CM 性：最高次局部上同调对应的环境 Ext 总能给出典范模。CM 性所增加的是其他 Ext 次数全部消失，以及典范模的深度达到 \(d\)。环境对偶信息因此集中在一个次数，下一小节才能把全部局部对偶写成只含 \(\omega_R\) 的公式。

\subsection{局部上同调与 Ext 的 Matlis 对偶}
\label{b3:subsec:2202:06}

为了从环境环 \(S\) 转到商环 \(R\)，需要把两个 Ext 计算联系起来。若 \(R\) 相对于 \(S\) 的对偶信息集中在次数 \(c\)，先通过 \(R\) 计算任意 \(R\) 模的环境 Ext，就只产生整体的次数平移 \(c\)。下面的引理将把这项平移写成两个 Ext 的自然同构。

\begin{lemma}\label{b3:lem:canonical-change-of-rings}
设 \(S\) 为交换 Noether 环，\(I\subseteq S\) 为理想，\(R=S/I\)。假设 \(S\) 作为自身上的模有有界内射消解，且存在整数 \(c\ge0\) 使
\(\operatorname{Ext}_S^q(R,S)=0\) 对 \(q\ne c\) 成立。记 \(W=\operatorname{Ext}_S^c(R,S)\)。则对每个有限生成的 \(R\) 模 \(M\) 及每个 \(p\ge0\)，有自然同构
\[
\operatorname{Ext}_R^p(M,W)
\cong\operatorname{Ext}_S^{p+c}(M,S),
\]
且 \(\operatorname{Ext}_S^q(M,S)=0\) 对 \(q<c\) 成立。
\end{lemma}
\begin{proof}
取内射 \(S\) 消解 \(J^\bullet\)，选整数 \(h\ge0\) 使 \(J^q=0\) 对 \(q>h\) 成立，再取每项有限秩自由的 \(R\) 消解 \(P_\bullet\rightarrow M\)。后一个消解可以无限延伸。令
\[
T^q=\operatorname{Hom}_S(R,J^q),
\qquad B^{p,q}=\operatorname{Hom}_R(P_p,T^q).
\]
每个 \(T^q\) 都是内射 \(R\) 模，理由是
\(\operatorname{Hom}_R(-,T^q)\cong\operatorname{Hom}_S(-,J^q)\)
正合。\(B\) 的指标范围为 \(p\ge0\)、\(0\le q\le h\)，总次数为 \(p+q\)。在 \((p,q)\) 项上，总微分取来自 \(P\) 的 Hom 微分加 \((-1)^p\) 倍来自 \(T\) 的微分，使两方向反交换。第一象限与 \(q\) 方向有界保证每条总次数对角线有限。

两种计算的目标分别是原来的环境 Ext 和经过 \(R\) 的 Ext。先沿 \(p\) 方向取上同调，\(T^q\) 的内射性使 \(p>0\) 时消失，\(p=0\) 为 \(\operatorname{Hom}_R(M,T^q)\cong\operatorname{Hom}_S(M,J^q)\)，所以总上同调是 \(\operatorname{Ext}_S^*(M,S)\)。再沿 \(q\) 方向取上同调，\(P_p\) 自由使 Hom 保持核、像及其商，故
\[
H^q(B^{p,\bullet})
=\operatorname{Hom}_R(P_p,H^q(T^\bullet)),
\]
它只在 \(q=c\) 留下 \(\operatorname{Hom}_R(P_p,W)\)。因此总次数 \(p+c\) 的上同调就是 \(\operatorname{Ext}_R^p(M,W)\)；低于 \(c\) 的总次数没有这一行的项，上同调为零。这里的单行比较使用次数滤过：逐行提升圈，利用 \(q\ne c\) 时的正合性消去相应分量，剩余行上的边界恰为总边界的像。每个固定总次数及相邻次数都只有有限项，因此圈与边界的比较只经过有限步。此构造与同态诱导的消解映射相容，给出自然性。
\end{proof}

\begin{theorem}[局部对偶]\label{b3:thm:canonical-module-local-duality}
设 \((S,\mathfrak n,k)\) 为 \(n\) 维交换完备 Noether 正则局部环，\(I\subsetneq S\) 为理想，且 \((R,\mathfrak m,k)=(S/I,\mathfrak n/I,k)\) 为 \(d\) 维 Cohen–Macaulay 环。取
\[
\omega_R=\operatorname{Ext}_S^{n-d}(R,S),
\qquad D_R=\operatorname{Hom}_R(-,E_R(k)).
\]
则对每个有限生成的 \(R\) 模 \(M\) 和每个整数 \(i\)，有关于 \(M\) 自然的同构
\begin{equation}\label{b3:eq:canonical-local-duality}
\begin{aligned}
H^i_{\mathfrak m}(M)
&\cong D_R\bigl(\operatorname{Ext}_R^{d-i}(M,\omega_R)\bigr),\\
D_R\bigl(H^i_{\mathfrak m}(M)\bigr)
&\cong\operatorname{Ext}_R^{d-i}(M,\omega_R).
\end{aligned}
\end{equation}
负次局部上同调和负次 Ext 均取零。所有 \(H^i_{\mathfrak m}(M)\) 都是 Artin 模，并且
\[
\operatorname{Ext}_R^p(M,\omega_R)=0\qquad(p>d).
\]
\end{theorem}
\begin{proof}
令 \(c=n-d\)。\cref{b3:prop:canonical-module-quotient}给出 \(\operatorname{Ext}_S^q(R,S)\) 只在 \(q=c\) 非零。\(S\) 的有界内射消解已在\cref{b3:prop:regular-local-top-cohomology}的证明中构造。故前一引理给出
\[
\operatorname{Ext}_S^{n-i}(M,S)
\cong\operatorname{Ext}_R^{d-i}(M,\omega_R)
\]
在 \(d-i\ge0\) 时成立；若 \(d-i<0\)，左端也由引理的低次消失为零。再用\cref{b3:thm:regular-local-duality}，以及已验证的 \(D_S(M)=D_R(M)\)、\(H^i_{\mathfrak n}(M)=H^i_{\mathfrak m}(M)\)，得到两行公式。Artin 性来自正则局部对偶；当 \(i<0\) 时，局部上同调为零，第二行给出最后的 Ext 消失。
\end{proof}

上式称为\term[jubu-duiou]{局部对偶}[local duality]。次数变换 \(q=d-i\) 中的 \(d\) 是底环 \(R\) 的维数；模自身的维数则限制哪些局部上同调次数可能出现，两者不能互换。公式允许用有限生成模的消解研究通常不有限生成的局部上同调，但没有把后者断言为有限生成模。例如在 \(k[\![t]\!]\) 上，\(\omega_R\cong R\)，公式的最高次仍然是 \(H^1_{(t)}(R)\cong E_R(k)\)，与节首计算一致。

\ProseParagraphBreak
设 \(R=S/I\) 满足上述局部对偶定理的假设，即 \(S\) 是完备 Noether 正则局部环，\(R\) 是 \(d\) 维 Cohen–Macaulay 商环，\(\mathfrak m\) 是 \(R\) 的极大理想，\(\omega_R\) 是所得典范模。对非零有限生成 \(R\) 模 \(M\)，令 \(\alpha=\operatorname{depth}_R M\)、\(t=\dim_R M\)。Matlis 对偶检测非零模，因此次数变换 \(q=d-i\) 同时保留消失与非消失：
\[
H^i_{\mathfrak m}(M)\ne0
\quad\Longleftrightarrow\quad
\operatorname{Ext}_R^{d-i}(M,\omega_R)\ne0.
\]
\cref{b3:tab:local-duality-dictionary}把深度、维数与 CM 条件的表述集中在一起。其中 \(M\) 为 CM 模并不要求 \(t=d\)。

\begin{table}[htbp]
\centering
\caption[局部对偶中的消失区间]{局部对偶中的消失区间。\(R\) 为上述完备 CM 商环，\(M\ne0\) 有限生成，\(\alpha=\operatorname{depth}_R M\)、\(t=\dim_R M\)。}
\label{b3:tab:local-duality-dictionary}
\begin{tabularx}{\linewidth}{@{}p{0.15\linewidth}XX@{}}
\toprule
数学信息 & 局部上同调 \(H^i_{\mathfrak m}(M)\) & \(\operatorname{Ext}_R^q(M,\omega_R)\) \\
\midrule
深度 \(\alpha\) & \(i<\alpha\) 时为零，\(i=\alpha\) 时非零 & \(q>d-\alpha\) 时为零，\(q=d-\alpha\) 时非零 \\
维数 \(t\) & \(i>t\) 时为零 & \(q<d-t\) 时为零 \\
\(M\) 为 CM 模 & 只在 \(i=t\) 非零 & 只在 \(q=d-t\) 非零 \\
\bottomrule
\end{tabularx}
\end{table}

可能非零的局部上同调次数满足 \(\alpha\le i\le t\)，次数变换将其翻转为 \(d-t\le q\le d-\alpha\)。深度给出 Ext 的最高非零次数 \(d-\alpha\)；模维数给出最低可能次数 \(d-t\)，即其下方全部消失。在 CM 情形 \(\alpha=t\)，整个区间缩成一个确实非零的次数。

\ProseParagraphBreak
先看模维数与底环维数确实不同的情形。取域 \(k\)，令 \(S=R=k[\![x,y]\!]\)、\(\mathfrak m=(x,y)\)、\(M=R/(x)\)。此时环境环维数 \(n\) 与底环维数 \(d\) 均为二，而 \(t=\dim_R M=1\)、\(\alpha=\operatorname{depth}_R M=1\)；\(y\) 在 \(M\cong k[\![y]\!]\) 上是非零因子。\(R\) 的典范模为 \(R\)，自由消解
\[
0\longrightarrow R\xrightarrow{x}R\longrightarrow M\longrightarrow0
\]
对偶后给出 \(\operatorname{Hom}_R(M,R)=0\)、\(\operatorname{Ext}^1_R(M,R)=R/(x)\)，且更高 Ext 为零。因此局部对偶给出
\[
D_R\bigl(H^1_{\mathfrak m}(M)\bigr)
\cong\operatorname{Ext}^{2-1}_R(M,R)
\cong R/(x).
\]
翻转次数中的二来自底环维数 \(d\)，不是模维数 \(t\)；模维数一则决定这里的局部上同调只可能在一次非零。

\ProseParagraphBreak
最后用\subsecref{b3:subsec:2201:03}的例子检验 CM 假设的作用。取域 \(k\) 及
\[
S=k[\![x,y]\!],\qquad R=S/(x^2,xy),\qquad\mathfrak m=(\bar x,\bar y).
\]
此前已算得
\[
\dim R=1,
\qquad H^0_{\mathfrak m}(R)=k\bar x\ne0.
\]
这个环还满足 \(H^1_{\mathfrak m}(R)=k(\!(y)\!)/k[\![y]\!]\ne0\)。环境环上的正则局部对偶分别给出
\[
\begin{aligned}
D_R(H^1_{\mathfrak m}(R))&\cong\operatorname{Ext}_S^1(R,S),\\
D_R(H^0_{\mathfrak m}(R))&\cong\operatorname{Ext}_S^2(R,S).
\end{aligned}
\]
两项均非零。\cref{b3:prop:canonical-module-quotient}仍给出典范模 \(\omega_R=\operatorname{Ext}_S^1(R,S)\)，但它只保留与最高次局部上同调对应的一项。若把局部对偶的单模公式用于 \(M=R\)、\(i=0\)，就会得到
\[
0\ne D_R\bigl(H^0_{\mathfrak m}(R)\bigr)
\cong\operatorname{Ext}_R^1(R,\omega_R)=0,
\]
矛盾。真正保留这部分信息的是另一个非零模 \(\operatorname{Ext}_S^2(R,S)\cong D_R(H^0_{\mathfrak m}(R))\)。因此一般情形须保留具有多个上同调的对偶复形，不能只取最高次局部上同调对应的那一个模。

\subsection{Cohen 结构定理与正则局部环的满射表示}
\label{b3:subsec:2202:07}

本节所有存在性及局部对偶证明都从给定满射 \(S\twoheadrightarrow R\) 开始。对显式幂级数商环，这份数据已经在题目中；例如 \(k[\![x_1,\ldots,x_n]\!]/I\) 的环境环就是所需的完备正则局部环。混合特征时也可以从给定完备离散赋值环 \(V\) 上的幂级数环 \(V[\![x_1,\ldots,x_r]\!]\) 开始，而不把系数环误写成域。

\ProseParagraphBreak
对一个未给出满射表示的交换完备 Noether 局部环，附录\cref{b3:thm:coefficient-field-existence,b3:thm:equicharacteristic-power-series-cover,b3:thm:mixed-characteristic-coefficient-ring}给出所需的环境环与满射。等特征时先取系数域，再按极大理想阶数展开，得到幂级数环满射；混合特征时以 Cohen 环替代系数域，同样逐阶选择剩余系数并取完备极限，得到其上幂级数环的满射。两种环境环均完备正则局部。相应结构定理还可对照松村英之\cite[\S 29, Theorems 29.3--29.4, pp. 225--226]{Matsumura1989CommutativeRingTheory}。

\ProseParagraphBreak
因此每个交换完备 Noether 局部环都有本节意义下的典范模；若该环还为 CM 环，就可使用\eqref{b3:eq:canonical-local-duality}。对于非完备局部环，这一构造首先给出的是完备化上的典范模；能否下降为原环上的有限生成模，是另一个需要条件的问题。

% !TeX root = ../../Book3.tex
% 纲要标题：### 22.3 Gorenstein 环
% 纲要位置：计划纲要.md，第 2204 行。

\section{Gorenstein 环}
\label{b3:sec:2203}

环自身的有限内射消解，比所有模具有有限自由消解的要求宽松。一般 Bass 判别把这种有限性转为 Cohen–Macaulay 性与顶次 Ext 的一维条件，完全交则可由 Koszul 对偶满足这一条件。


% 纲要内容：
%
% * Noether 局部环的有限自内射维数
% * 与 Cohen–Macaulay 性的联系
% * 完全交蕴含 Gorenstein 的适当版本
%

\subsection{Noether 局部环的有限自内射维数}
\label{b3:subsec:2203:01}

正则局部环用有限自由消解控制所有模；Gorenstein 性则问环本身能否有有限的内射消解。两种要求的方向不同。特别是，剩余域有无限自由消解，并不排除环作为自身的模内射。

\begin{definition}\label{b3:def:injective-dimension-gorenstein}
设 \(R\) 为交换环，\(N\ne0\) 为 \(R\) 模。若存在正合列
\[
0\longrightarrow N\longrightarrow E^0\longrightarrow\cdots
\longrightarrow E^s\longrightarrow0
\]
且每个 \(E^i\) 内射，则满足此条件的最小 \(s\) 称为 \(N\) 的\term[neisheweishu]{内射维数}[injective dimension]，记为 \(\operatorname{id}_RN\)；若不存在有限长度的此种消解，则记为 \(\infty\)。对交换 Noether 局部环 \(R\)，条件 \(\operatorname{id}_RR<\infty\) 成立时，称 \(R\) 为 \term[Gorensteinhuan]{Gorenstein 环}[Gorenstein ring]。交换 Noether 环在每个素理想处的局部环均有此性质时，也称为 Gorenstein 环。
\end{definition}
\symidx{\(\operatorname{id}_RN\)：内射维数}

内射消解中的模通常不是有限模。譬如离散赋值环 \(R\) 的分式域 \(K\) 给出
\[
0\longrightarrow R\longrightarrow K\longrightarrow K/R\longrightarrow0.
\]
两个右端模都可除，由主理想整环上的 Baer 判据都内射；而 \(R\) 不是内射，因为均匀参数的倒数不在 \(R\) 中。因此 \(\operatorname{id}_RR=1\)。这里仍是用一个有限长度的消解描述一个一维环，却不能要求消解各项有限生成。

\begin{lemma}\label{b3:lem:injective-dimension-test}
设 \(R\) 为交换环，\(N\ne0\) 为 \(R\) 模，\(s\ge0\)。下列两条件等价：
\[
\operatorname{id}_RN\le s;\qquad
\operatorname{Ext}_R^{s+1}(R/J,N)=0
\text{ 对每个理想 }J\subseteq R\text{ 成立}.
\]
\end{lemma}
\begin{proof}
第一条件推出更高次的 Ext 全部为零。反过来，在一个内射消解中截取
\(0\to N\to E^0\to\cdots\to E^{s-1}\to C\to0\)；\(s=0\) 时令 \(C=N\)。逐次使用第二变量长正合列，得到
\(\operatorname{Ext}_R^1(R/J,C)\cong\operatorname{Ext}_R^{s+1}(R/J,N)=0\)。
由 \(0\to J\to R\to R/J\to0\) 的 Hom 长正合列，每个 \(J\to C\) 均可延拓到 \(R\)。第二卷第7.2节的 Baer 判据使 \(C\) 内射，故上述截断就是长度至多 \(s\) 的内射消解。

这里使用的内射计算与此前用自由消解定义的 Ext 一致：把自由消解 \(P_\bullet\to R/J\) 和内射消解 \(N\to E^\bullet\) 组成 \(\operatorname{Hom}_R(P_p,E^q)\)，两个方向分别由自由性和内射性保持正合，因而分别只剩 \(\operatorname{Hom}_R(P_\bullet,N)\) 与 \(\operatorname{Hom}_R(R/J,E^\bullet)\)。给定总次数中只有有限个项，双复形的圈消去论证给出相同上同调。
\end{proof}

\subsection{Gorenstein 性与 Cohen–Macaulay 性}
\label{b3:subsec:2203:02}

有限自内射维数不仅给出一项同调界，还强迫环达到 Cohen--Macaulay 深度。下面的判别不要求环完备，也不要求事先给定正则局部覆盖。证明采用素点间的 Ext 比较与正则元素换环，可对照松村英之\cite[\S 18, Lemmas 1--4 and Theorem 18.1, pp. 139--144]{Matsumura1989CommutativeRingTheory}。

\begin{theorem}[Gorenstein 环的判别]\label{b3:thm:gorenstein-bass-criterion}
设 \((R,\mathfrak m,k)\) 为交换 Noether 局部环，\(d=\dim R\)。下列条件等价：
\begin{equivalences}
\item \(R\) 为 Gorenstein 环；
\item \(R\) 为 Cohen--Macaulay 环，且
\(\dim_k\operatorname{Ext}_R^d(k,R)=1\)；
\item \(\operatorname{Ext}_R^i(k,R)=0\) 对 \(i\ne d\) 成立，且
\(\operatorname{Ext}_R^d(k,R)\cong k\)。
\end{equivalences}
这些条件成立时，\(\operatorname{id}_RR=d\)。
\end{theorem}
\begin{proof}
先说明三个计算事实。第一，\cref{b3:lem:injective-dimension-test}中的理想可只取素理想：每个有限模有有限素循环滤过，长正合列把素循环商上的 Ext 消失传到任意有限模，特别是各个 \(R/J\)。

第二，若 \(x\in\mathfrak m\) 是 \(R\) 非零因子，令 \(\overline R=R/xR\)。对每个 \(\overline R\) 模 \(N\)，有
\[
\operatorname{Hom}_R(N,R)=0,\qquad
\operatorname{Ext}_R^{j+1}(N,R)
\cong\operatorname{Ext}_{\overline R}^{j}(N,\overline R)
\quad(j\ge0).
\]
确实，二项自由消解 \(0\to R\xrightarrow{x}R\to\overline R\to0\) 给出
\(\operatorname{Ext}_R^q(\overline R,R)\) 只在 \(q=1\) 为 \(\overline R\)。把 \(N\) 的自由 \(\overline R\) 消解与 \(R\) 的内射消解组成 Hom 双复形，按\cref{b3:lem:canonical-change-of-rings}证明中的单行消去，便得上述换环同构；每个总次数只有有限项，不涉及无限乘积。

第三，对有限 \(R\) 模 \(M\)、素理想 \(\mathfrak p\) 及 \(t=\dim R/\mathfrak p\)，有
\[
\operatorname{Ext}_R^{j+t}(k,M)=0
\ \Longrightarrow\
\operatorname{Ext}_{R_{\mathfrak p}}^j
  (\kappa(\mathfrak p),M_{\mathfrak p})=0.
\]
先设 \(t=1\)，取 \(a\in\mathfrak m\setminus\mathfrak p\)。模 \(R/(\mathfrak p,a)\) 有有限长度，故给定的 \((j+1)\) 次消失也对它成立。由
\(0\to R/\mathfrak p\xrightarrow{a}R/\mathfrak p\to R/(\mathfrak p,a)\to0\)，乘以 \(a\) 在有限模 \(\operatorname{Ext}_R^j(R/\mathfrak p,M)\) 上满射；Nakayama 引理使该模为零，局部化即得断言。一般 \(t\) 时，取长度为 \(t\) 的最长素理想链，从 \(\mathfrak m\) 到 \(\mathfrak p\) 逐个作这一步；每个相邻商的局部维数为一，Ext 与有限模的局部化相容。

现在设 \(r=\operatorname{id}_RR<\infty\)。取满足 \(\dim R/\mathfrak p=d\) 的极小素理想，则零维局部环 \(R_{\mathfrak p}\) 有非零基座。第三个事实的逆否命题给出 \(\operatorname{Ext}_R^d(k,R)\ne0\)，所以 \(r\ge d\)。若 \(r>0\)，则 \(\operatorname{Ext}_R^r(k,R)\ne0\)：第一事实给出某个 \(\mathfrak p\) 使 \(\operatorname{Ext}_R^r(R/\mathfrak p,R)\ne0\)；若 \(\mathfrak p\ne\mathfrak m\)，取 \(a\in\mathfrak m\setminus\mathfrak p\)，同一短正合列及 \((r+1)\) 次 Ext 消失会使乘以 \(a\) 满射，再由 Nakayama 引理得到矛盾。

此时 \(\mathfrak m\notin\operatorname{Ass}R\)。否则有单射 \(k\to R\)，其余核的 \((r+1)\) 次 Ext 为零，迫使
\(0=\operatorname{Ext}_R^r(R,R)\to\operatorname{Ext}_R^r(k,R)\) 满射，矛盾。因此可选 \(x\in\mathfrak m\) 为非零因子。第二事实和第一事实使 \(\operatorname{id}_{\overline R}\overline R=r-1\)：上界由全部素循环商的换环同构得到，下界由 \(\operatorname{Ext}_R^r(k,R)\ne0\) 得到。对 \(r\) 归纳，\(r=0\) 时已有 \(d=0\)，得到 \(R\) 为 CM 环且 \(r=d\)。

归纳还给出条件3。零维且自内射时，任取单射 \(k\to R\)，内射性使 \(\operatorname{Hom}_R(R,R)\to\operatorname{Hom}_R(k,R)\) 满射。后者由该单射生成且被 \(\mathfrak m\) 零化，故同构于 \(k\)，正次 Ext 全为零。正维时，\(\operatorname{Hom}_R(k,R)=0\)，其余各次由第二事实下降到 \(\overline R\)。这证明条件1推出条件3；条件3结合深度的 Ext 刻画显然推出条件2。

最后设条件2成立。取参数正则序列 \(x_1,\ldots,x_d\)，令 \(T=R/(x_1,\ldots,x_d)\)。反复换环得到
\(\operatorname{Soc}(T)\cong\operatorname{Ext}_R^d(k,R)\cong k\)。Artin 局部环中每个非零理想都与基座非零相交：对该理想反复乘以极大理想，取最后一个非零幂次即可。将基座同构 \(\operatorname{Soc}(T)\to k\subseteq E_T(k)\) 延拓到 \(T\to E_T(k)\)，本质性使其为单射；由\cref{b3:lem:matlis-finite-length}，
\(\ell_T E_T(k)=\ell_T D_T(T)=\ell_T T\)，故它是同构，\(T\) 自内射。

沿正则序列逐项返回。若 \(\operatorname{id}_{\overline R}\overline R\le s\)，则对含 \(x\) 的素理想，换环给出 \((s+2)\) 次 Ext 消失；对不含 \(x\) 的素理想，模 \(R/(\mathfrak p,x)\) 被 \(x\) 零化，短正合列使乘以 \(x\) 在 \(\operatorname{Ext}_R^{s+2}(R/\mathfrak p,R)\) 上满射，Nakayama 引理再次使其为零。第一事实于是给出 \(\operatorname{id}_RR\le s+1\)。从 \(T\) 返回 \(R\) 得 \(\operatorname{id}_RR\le d\)，而顶次 Ext 非零给出等号，完成证明。
\end{proof}

对一个 \(d\) 维 CM 局部环 \(R\)，整数
\(\dim_k\operatorname{Ext}_R^d(k,R)\) 称为它的\term[CohenMacaulayleixing]{Cohen–Macaulay 类型}[Cohen--Macaulay type]。深度的 Ext 刻画保证这个整数为正。判别定理说明 Gorenstein 环恰是类型为一的 CM 环；只有“CM”仍不足够。零维环的类型就是其基座的维数，所以 \(k[x,y]/(x,y)^2\) 是 CM 环但不是 Gorenstein 环。

\begin{proposition}\label{b3:prop:gorenstein-completion-localization}
设 \((R,\mathfrak m,k)\) 为交换 Noether 局部环。则 \(R\) 为 Gorenstein 环当且仅当其 \(\mathfrak m\)-进完备化 \(\widehat R\) 为 Gorenstein 环。若 \(R\) 为 Gorenstein 环，则每个 \(R_{\mathfrak p}\)、\(\mathfrak p\in\operatorname{Spec}R\)，也为 Gorenstein 环。
\end{proposition}
\begin{proof}
完备化保维数并忠实平坦，且由\cref{b3:thm:cm-completion}保持及反映 CM 性。取 \(k\) 的有限秩自由 \(R\) 消解，逐项张量 \(\widehat R\) 后仍是 \(k\) 的自由消解。有限秩保证 Hom 与此张量相容，平坦性保证上同调与张量相容，故
\[
\operatorname{Ext}_R^i(k,R)\otimes_R\widehat R
\cong\operatorname{Ext}_{\widehat R}^i(k,\widehat R).
\]
左侧原本被 \(\mathfrak m\) 零化，张量后不变。应用上一判别定理，便得完备化等价性。

为处理局部化，先说明内射模 \(E\) 的局部化 \(E_{\mathfrak p}\) 对 \(R_{\mathfrak p}\) 内射。该局部环的每个理想都是某个有限生成理想 \(J\subseteq R\) 的局部化。Noether 性使 \(J\) 有限展示，从而
\[
\operatorname{Hom}_{R_{\mathfrak p}}(J_{\mathfrak p},E_{\mathfrak p})
\cong\operatorname{Hom}_R(J,E)_{\mathfrak p}.
\]
原内射性使 \(\operatorname{Hom}_R(R,E)\to\operatorname{Hom}_R(J,E)\) 满射，局部化后仍满射，再用 Baer 判据即可。把 \(R\) 的有限内射消解逐项局部化，就得到 \(R_{\mathfrak p}\) 的有限内射消解。
\end{proof}

在给定正则覆盖的情形，典范模提供另一种直接的充分条件，并说明自内射维数为何能由局部对偶控制。

\begin{proposition}\label{b3:prop:canonical-free-gorenstein}
设 \((S,\mathfrak n)\) 为 \(n\) 维交换 Noether 完备正则局部环，\(I\subsetneq S\) 为理想，\(R=S/I\) 为 \(d\) 维 Cohen--Macaulay 环，典范模为
\(\omega_R=\operatorname{Ext}_S^{n-d}(R,S)\)。若 \(\omega_R\cong R\) 为 \(R\) 模，则 \(R\) 为 Gorenstein 环。
\end{proposition}
\begin{proof}
对每个理想 \(J\subseteq R\)，模 \(R/J\) 有限。\cref{b3:thm:canonical-module-local-duality}给出
\[
\operatorname{Ext}_R^{d+1}(R/J,\omega_R)
\cong D_R\bigl(H_{\mathfrak m}^{-1}(R/J)\bigr)=0.
\]
用给定的 \(\omega_R\cong R\) 替换第二变量，再由\cref{b3:lem:injective-dimension-test}得 \(\operatorname{id}_RR\le d\)。此外 CM 性及\cref{b3:thm:depth-ext}使 \(\operatorname{Ext}_R^d(k,R)\ne0\)，所以内射维数恰为 \(d\)。
\end{proof}

此处典范模同构是一个 \(R\) 模同构，通常没有指定的自然选择。把典范模选成 \(R\) 也不等于对每个有限模使用普通线性对偶 \(\operatorname{Hom}_R(-,R)\)：局部对偶中的次数 \(d-i\) 和 Matlis 对偶仍然不可省略。

\subsection{完全交的 Gorenstein 性}
\label{b3:subsec:2203:03}

完全交商的 Koszul 消解有一种对称性：取 \(S\) 线性对偶后，复形各项的顺序倒过来，微分仍是原正则序列给出的矩阵。这使典范模的计算十分具体。

\begin{proposition}\label{b3:prop:ci-canonical-module}
设 \((S,\mathfrak n)\) 为 \(n\) 维交换 Noether 完备正则局部环，\(f_1,\ldots,f_c\in\mathfrak n\) 为 \(S\) 正则序列，\(R=S/(f_1,\ldots,f_c)\)。则 \(d=\dim R=n-c\)，而且
\[
\operatorname{Ext}_S^j(R,S)=
\begin{cases}
R,&j=c,\\
0,&j\ne c.
\end{cases}
\]
特别地，\(R\) 的典范模 \(\omega_R\) 同构于 \(R\)。
\end{proposition}
\begin{proof}
\cref{b3:thm:complete-intersection-cm}及\cref{b3:prop:given-complete-intersection}给出 CM 性和维数。用\cref{b3:prop:ci-koszul-resolution}的 Koszul 消解 \(K_\bullet(f;S)\) 计算 Ext。外积给出完美配对
\[
\bigwedge^jS^c\times\bigwedge^{c-j}S^c
\longrightarrow\bigwedge^cS^c\cong S,
\]
最后一个同构取有序基 \(e_1\wedge\cdots\wedge e_c\mapsto1\)。若 \(u\) 的次数为 \(j\)，Koszul 微分满足
\(\mathrm{d}(u\wedge v)=\mathrm{d}u\wedge v+(-1)^ju\wedge \mathrm{d}v\)。因此配对使对偶微分与倒序 Koszul 微分相差一个只依赖次数的符号；逐次把某些次数上的配对乘以 \(-1\)，便得到复形同构
\[
\operatorname{Hom}_S(K_\bullet(f;S),S)^j
\cong K_{c-j}(f;S).
\]
正则序列使右侧只有链次数零的同调 \(R\)，所以对偶上同调只在 \(j=c\) 为 \(R\)。\(c=0\) 时消解只有 \(S\)，结论同样成立。由于 \(n-d=c\)，这正是典范模定义要求的次数。
\end{proof}

\begin{theorem}\label{b3:thm:complete-intersection-gorenstein}
每个交换 Noether 完全交局部环都是 Gorenstein 环。因此，对交换 Noether 局部环有
\[
\text{正则}\ \Longrightarrow\ \text{完全交}\
\Longrightarrow\ \text{Gorenstein}\
\Longrightarrow\ \text{Cohen--Macaulay}.
\]
\end{theorem}
\begin{proof}
按照\cref{b3:def:local-complete-intersection}，完全交局部环 \(R\) 的完备化具有
\(\widehat R=S/(f_1,\ldots,f_c)\) 的表示，其中 \(S\) 完备正则局部，\(f_1,\ldots,f_c\) 为正则序列。前一命题使 \(\omega_{\widehat R}\cong\widehat R\)，再由\cref{b3:prop:canonical-free-gorenstein}得 \(\widehat R\) 为 Gorenstein 环。\cref{b3:prop:gorenstein-completion-localization}把此性质下降到 \(R\)。链中的第一项由\cref{b3:thm:complete-intersection-cm}证明中的空正则序列论证成立，最后一项来自\cref{b3:thm:gorenstein-bass-criterion}。
\end{proof}

完全交的定义已给出 \(\widehat R\) 的正则局部覆盖及其正则序列关系，所以此处可直接计算典范模。对一般交换完备 Noether 局部环，覆盖的存在性由\appref{b3:app:D}的 Cohen 结构定理保证，但其关系未必由正则序列生成。对于给定的正则局部环 \(S\) 及非零非单位 \(f\)，商 \(S/(f)\) 是超曲面，因而也是 Gorenstein 环；它是否正则仍须另外判断。

% !TeX root = ../../Book3.tex
% 纲要标题：### 22.4 Gorenstein 环的对偶例子
% 纲要位置：计划纲要.md，第 2692 行。

\section{Gorenstein 环的对偶例子}
\label{b3:sec:2204}

零维 Gorenstein 性可以通过基座及乘法配对检验，因而适合具体计算。正维及非 Cohen–Macaulay 情形还要保留对偶的次数；一个典范模未必足以记录全部数据。


% 纲要内容（仅作写作计划，尚非教材正文）：
% * Artin Gorenstein 环的基座与乘法配对
% * 超曲面及完全交的环境环 Ext 计算
% * 环境环上的对偶复形：有限长度和各项有限秩的自由消解、对偶微分及移位；区分各项的 $S$ 模结构与上同调的 $R$ 模结构
% * 非 CM 例子的两个 Ext 次数与局部上同调的 Matlis 对应
% ---

\subsection{Artin Gorenstein 环}
\label{b3:subsec:2204:01}

零维时，Gorenstein 性可以通过一个很小的线性空间辨认。基座包含那些被极大理想全部零化的元素；它们是乘法再也不能向下推进的部分。

\begin{proposition}\label{b3:prop:artinian-gorenstein-socle}
设 \((R,\mathfrak m,k)\) 为非零交换 Artin 局部环，并记 \(\operatorname{Soc}(R)=(0:_R\mathfrak m)\)。则下列条件等价：\(R\) 为 Gorenstein 环；\(R\) 为内射 \(R\) 模；\(\dim_k\operatorname{Soc}(R)=1\)。
\end{proposition}
\begin{proof}
零维时，\cref{b3:thm:gorenstein-bass-criterion}给出内射维数为零，并使
\(\operatorname{Hom}_R(k,R)=\operatorname{Soc}(R)\) 为一维。下面也直接说明基座条件如何产生自内射性。Artin 局部环完备，其有限长度模 \(R\) 的 Matlis 对偶 \(D_R(R)=E_R(k)\) 与 \(R\) 有相同长度；这是\cref{b3:thm:matlis-duality}在组成列上的应用，因为 \(D_R(k)\cong k\) 且对偶正合。

若基座一维，把 \(\operatorname{Soc}(R)\cong k\) 嵌入 \(E_R(k)\)，利用内射性将它延拓为 \(u:R\to E_R(k)\)。任一非零理想 \(J\) 都与基座非零相交：因 \(\mathfrak m\) 幂零，选最大 \(r\) 使 \(\mathfrak m^rJ\ne0\)，则该模被 \(\mathfrak m\) 零化。因此 \(\ker u\) 若非零，必与基座相交，与 \(u\) 在基座上单射矛盾。故 \(u\) 单射，两端长度相同使它为同构，\(R\) 遂内射。
\end{proof}

若 \(R=A\) 还是剩余域为 \(k\) 的有限维交换局部 \(k\) 代数，上述结论可直接写成乘法配对。选 \(k\) 线性泛函 \(\lambda:A\to k\)，使它在一维基座上非零。则
\[
A\longrightarrow\operatorname{Hom}_k(A,k),\qquad
a\longmapsto\bigl(b\longmapsto\lambda(ab)\bigr)
\]
单射：若 \(a\ne0\)，理想 \(Aa\) 包含非零基座元素，所以存在 \(b\) 使 \(\lambda(ab)\ne0\)。两端维数相同，故这是同构。这里的目标模取作用 \((c\phi)(b)=\phi(cb)\)；它内射，因为
\[
\operatorname{Hom}_A(N,\operatorname{Hom}_k(A,k))\cong\operatorname{Hom}_k(N,k).
\]
同构把 \(f\) 送到 \(n\mapsto f(n)(1)\)，逆映射把 \(u\) 送到 \(n\mapsto[a\mapsto u(an)]\)。向量空间对偶正合，从而此 Hom 函子正合。

\ProseParagraphBreak
例如 \(A=k[x,y]/(x^a,y^b)\)、\(a,b\ge1\)，基座由 \(x^{a-1}y^{b-1}\) 生成。取 \(\lambda\) 为此单项式的系数，基单项式 \(x^iy^j\) 与 \(x^{a-1-i}y^{b-1-j}\) 配对为一，其余达到不同次数的乘积对该系数贡献为零。这个有限维配对把“类型一”的条件变成能够逐项检验的矩阵。同样长度很短的 \(k[x,y]/(x,y)^2\) 却有二维基座 \(kx+ky\)，所以不存在这样的非退化乘法配对。

\ProseParagraphBreak
对后一环 \(A=k[x,y]/(x,y)^2\)，典范模仍然存在，但不再同构于环本身。令 \(W=\operatorname{Hom}_k(A,k)\)，以 \(\phi_0,\phi_x,\phi_y\) 表示基 \(1,x,y\) 的对偶基。按照 \((a\phi)(b)=\phi(ab)\) 的作用，直接计算得
\[
\begin{aligned}
x\phi_x&=\phi_0,& y\phi_y&=\phi_0,\\
x\phi_y&=0,& y\phi_x&=0,\\
x\phi_0&=0,& y\phi_0&=0.
\end{aligned}
\]
于是 \(\operatorname{Soc}_A(W)=k\phi_0\)。任意非零子模中，若一个元素的 \(\phi_x\) 或 \(\phi_y\) 系数非零，乘 \(x\) 或 \(y\) 就得到非零的 \(\phi_0\)；否则它本身已在 \(k\phi_0\) 中。因此 \(k\phi_0\subseteq W\) 本质。上面的 Hom 计算又说明 \(W\) 内射，故 \(W\cong E_A(k)=D_A(A)\)。零维时 \(H^0_{\mathfrak m}(A)=A\)，Matlis 双对偶遂使 \(W\) 成为典范模。可是 \(\mathfrak mW=k\phi_0\)，所以 \(W/\mathfrak mW\) 为二维，\(W\) 至少需要两个生成元；\(A\) 作为自身模只需一个。这个作用表具体说明了 \(\omega_A\not\cong A\)，也区分了典范模的存在与 Gorenstein 性。

\subsection{超曲面与完全交的对偶例子}
\label{b3:subsec:2204:02}

设 \(k\) 为域，取 \(S=k[\![x,y]\!]\)、\(R=S/(xy)\)。环境环中的自由消解是
\[
0\longrightarrow S\xrightarrow{xy}S\longrightarrow R\longrightarrow0.
\]
对偶后仍为乘 \(xy\)，所以 \(\omega_R=\operatorname{Ext}_S^1(R,S)\cong R\)。这是一维 Gorenstein 环，但在原点有两条分支，因而不是整环，也不是正则局部环。Gorenstein 条件描述对偶性质，不能替代不可约性或光滑性。

\ProseParagraphBreak
该例的局部上同调也可以算出。因为在 \(R\) 中 \(xy=0\)，有 \(R_{xy}=0\)，所以两元 Čech 复形给出
\[
H_{\mathfrak m}^1(R)=\bigl(R_x\oplus R_y\bigr)/R
=\bigl(k((x))\oplus k((y))\bigr)/R.
\]
此处 \(R\) 在两支中的像是所有常数项相等的幂级数对：任一元素唯一写成
\(c+x\cdot a(x)+y\cdot b(y)\)。故上式中的 \(R\to R_x\oplus R_y\) 为单射，\(H_{\mathfrak m}^0(R)=0\)。局部对偶进一步给出
\(D_R(H_{\mathfrak m}^1(R))\cong R\)，并由 Matlis 双对偶得到
\(H_{\mathfrak m}^1(R)\cong E_R(k)\)。最后这个同构反映的是内射包络结构，不应被误写成 \(H_{\mathfrak m}^1(R)\cong R\)。

\ProseParagraphBreak
再设 \(k\) 为域，取 \(S=k[\![x,y,z]\!]\)、\(R=S/(x^2,y^3)\)。以 \((x^2,y^3)\) 为序，消解为
\[
0\longrightarrow S\xrightarrow{\binom{-y^3}{x^2}}S^2
\xrightarrow{(x^2\ \ y^3)}S\longrightarrow R\longrightarrow0.
\]
取对偶后，零次核为零；若 \(-y^3a+x^2b=0\)，则互素性给出
\((a,b)=(x^2c,y^3c)\)，故一次上同调为零；二次余核为 \(R\)。所以 \(\omega_R\cong R\)，且 \(\dim R=1\)。参数 \(z\) 为非零因子，商去 \(z\) 后得到前面有非退化配对的 Artin 环 \(k[x,y]/(x^2,y^3)\)。

\ProseParagraphBreak
这些对偶计算全部在环境环 \(S\) 上进行。商环 \(R\) 上的剩余域自由消解通常无限长；例如 \(k[\varepsilon]/(\varepsilon^2)\) 上由乘 \(\varepsilon\) 构成的无限消解已在\secref{b3:sec:1903}计算，而环本身却内射。必须在记号中保留 \(\operatorname{Ext}_S\)、\(\operatorname{Ext}_R\) 的下标，才能区分这两项事实。

\subsection{环境环上的对偶复形}
\label{b3:subsec:2204:03}

对完备正则局部环 \(S\) 的一个 CM 商 \(R=S/I\)，自由消解的对偶只有一个非零上同调模，因而典范模能够集中表达对偶信息。若 \(R\) 不是 CM 环，不同次数上的对偶信息可能同时存在；这时保留整个对偶复形，比只取最高次局部上同调的对偶更合适。

\begin{definition}\label{b3:def:regular-cover-dual-complex}
设 \((S,\mathfrak n,k)\) 为 \(n\) 维交换完备 Noether 正则局部环，\(I\subsetneq S\) 为理想，\(R=S/I\)。取有限长度、各项均为有限秩自由 \(S\) 模的消解 \(P_\bullet\rightarrow R\)，微分记为 \(\mathrm d_j:P_j\rightarrow P_{j-1}\)。令
\(C^j=\operatorname{Hom}_S(P_j,S)\)，微分为
\(\delta^j(\varphi)=\varphi\circ\mathrm d_{j+1}\)。把
\[
\operatorname{Hom}_S(P_\bullet,S)[n]
\]
看作上链复形，称它为给定满射表示 \(S\twoheadrightarrow R\) 下的一个归一化对偶复形模型。移位约定为
\(C[n]^j=C^{j+n}\)、\(\mathrm{d}_{C[n]}=(-1)^n\delta\)，所以对每个整数 \(i\)，该模型的上同调为
\[
H^{-i}(C[n])=\operatorname{Ext}_S^{n-i}(R,S).
\]
未出现的 \(P_j\) 及负次 Ext 均记为零。
\end{definition}

上述有限消解可以取得：\cref{b3:thm:regular-local-global-dimension}给出正则局部环 \(S\) 的整体维数为 \(n\)，而有限生成模的极小自由消解在每一项只有有限个生成元。若 \(P_\bullet\) 与 \(Q_\bullet\) 都是这样的消解，恒等映射 \(R\rightarrow R\) 可提升为两个方向的比较映射；其复合与各自的恒等映射链同伦。取 \(S\) 对偶并移位后，仍得到互为逆的上链同伦等价。因此这里的消解选择不改变模型的链同伦类型，也不改变各次上同调。

\ProseParagraphBreak
该模型的各项是 \(S\) 模，未必是 \(R\) 模，因为 \(I\) 未必零化这些自由项。不过，对任意 \(a\in I\)，消解上的乘 \(a\) 映射诱导 \(R\) 上的零映射，比较映射的同伦唯一性使它链同伦于零。对偶后乘 \(a\) 同样同伦于零，故 \(I\) 零化每个上同调模。这才给出上述 Ext 模的 \(R\) 模结构。

\ProseParagraphBreak
若 \(R\) 为 \(d\) 维 CM 环，\cref{b3:thm:regular-local-duality}及局部上同调的消失说明只有
\(\operatorname{Ext}_S^{n-d}(R,S)=\omega_R\) 非零。因此归一化后的模型在次数 \(-d\) 的上同调为 \(\omega_R\)，其余次数均为零；这个次数与局部对偶公式中的 \(d-i\) 相对应。

\ProseParagraphBreak
最后设 \(k\) 为域，取 \(S=k[\![x,y]\!]\)、\(R=S/(x^2,xy)\)，并记 \(\mathfrak m=(x,y)R\)。若 \(x^2a+xyb=0\)，在整环 \(S\) 中消去 \(x\) 得 \(xa+yb=0\)；由于 \(x,y\) 互素，必有 \((a,b)=c(-y,x)\)。故有自由消解
\[
0\longrightarrow S\xrightarrow{\binom{-y}{x}}S^2
\xrightarrow{(x^2\ \ xy)}S\longrightarrow R\longrightarrow0.
\]
取 \(S\) 对偶后，两个微分分别是
\[
S\xrightarrow{\binom{x^2}{xy}}S^2
\xrightarrow{(-y\ \ x)}S.
\]
若 \(-ya+xb=0\)，同样的整除论证给出 \((a,b)=c(x,y)\)。所以一次核由 \((x,y)\) 生成，而前一像由 \(x(x,y)\) 生成；二次余核则是 \(S/(x,y)\)。因此
\[
\operatorname{Ext}_S^1(R,S)\cong S/(x),\qquad
\operatorname{Ext}_S^2(R,S)\cong S/(x,y)=k.
\]
这里 \(S\) 的维数为二，归一化对偶模型在次数 \(-1\) 和零都有非零上同调。\subsecref{b3:subsec:2201:03}给出的两个局部上同调模正好与之对应。具体地，\(R\) 为完备 Noether 局部环，取
\(D_R(-)=\operatorname{Hom}_R(-,E_R(k))\)。由\cref{b3:thm:regular-local-duality}及商环上的对偶识别，得到\cref{b3:tab:non-cm-dual-degrees}中的两项同构；Ext 的下标仍是环境环 \(S\)。

\begin{table}[htbp]
\centering
\caption{非 Cohen--Macaulay 环的两个对偶次数}
\label{b3:tab:non-cm-dual-degrees}
\begin{tabularx}{\linewidth}{@{}cXX@{}}
\toprule
次数 \(i\) & 局部上同调 \(H^i_{\mathfrak m}(R)\) & Matlis 对偶 \(D_R(H^i_{\mathfrak m}(R))\) \\
\midrule
\(1\) & \(k((y))/k[\![y]\!]\) & \(\operatorname{Ext}_S^1(R,S)\cong S/(x)\) \\
\(0\) & \(kx\cong k\) & \(\operatorname{Ext}_S^2(R,S)\cong k\) \\
\bottomrule
\end{tabularx}
\end{table}

典范模 \(\omega_R=\operatorname{Ext}_S^1(R,S)\) 记录了一次局部上同调的对偶，却没有记录由 \(kx\) 产生的另一个非零次数。这个例子说明，非 CM 环的局部对偶信息通常需要多个模及其次数，不能仅由一个典范模替代。


\begin{chaptersummary}
局部上同调来自一个左正合函子的导出，但在 Noether 环上可由有限的局部化复形计算。局部环上有限生成模的深度给出首次非零次数，维数给出最高可能次数；CM 性恰好把全部信息集中到一个次数。Matlis 对偶使这些通常不有限生成的 Artin 模，能够与有限生成模相互转换，完备性保证双对偶返回原模。

\ProseParagraphBreak
给定从完备正则局部环到所研究局部环的满射后，环境环上的 Ext 构造典范模；CM 假设使其余次数消失，从而得到含单个典范模的局部对偶。完全交的 Koszul 自对偶给出 \(\omega_R\cong R\)，继而推出有限自内射维数。零维时，这一条件又成为基座一维和非退化乘法配对。非 CM 例子保留多个非零 Ext，说明一般对偶理论为何必须研究复形。
\end{chaptersummary}

\BookExercises

\begin{optionalexercise}\label{b3:ex:22-unbounded-torsion}
设 \(k\) 为域。在 \(R=k[x]\) 上令 \(M=\bigoplus_{n\ge1}R/(x^n)\)。证明 \(\Gamma_{(x)}(M)=M\)，但不存在统一的 \(N\) 使 \(x^NM=0\)。再计算全部 \(H_{(x)}^i(M)\)。
\end{optionalexercise}
\begin{solution}
每个元素只有有限个非零分量，取这些分量指标的最大值即可零化它。对于任意 \(N\)，第 \(N+1\) 个分量中的一不被 \(x^N\) 零化，所以没有统一指数。局部化 \(M_x=0\)，两项 Čech 复形为 \(M\to0\)，故零次为 \(M\)，所有正次数为零。
\end{solution}

\begin{optionalexercise}\label{b3:ex:22-principal-nonreduced}
设 \(k\) 为域。令 \(R=k[x,y]/(xy)\)，\(I=(x)\)。求 \(H_I^0(R)\)、\(H_I^1(R)\)，并说明它们分别记录了哪一部分元素。
\end{optionalexercise}
\begin{solution}
局部化后 \(x\) 可逆迫使 \(y=0\)，所以 \(R_x=k[x,x^{-1}]\)。映射 \(R\to R_x\) 的核为 \((y)\)，像为 \(k[x]\)。故
\[
H_I^0(R)=(y),\qquad
H_I^1(R)=k[x,x^{-1}]/k[x].
\]
第一项记录 \(y\) 轴这一不可约分量上被 \(x\) 零化的元素；第二项由负次幂的类组成，记录局部化后新增、却不能由原环提升的元素。两项性质不同，不能把局部上同调全都看成原模的子模。
\end{solution}

\begin{optionalexercise}\label{b3:ex:22-two-variable-cech}
设 \(k\) 为域，\(R=k[x,y]\)、\(\mathfrak m=(x,y)\)、\(M=R/(xy)\)。用二元 Čech 复形计算全部 \(H_{\mathfrak m}^i(M)\)，并给出每个非零模的 \(k\) 基。说明为什么一次上同调除两个分支的负次幂之外，还含有一份常数差。
\end{optionalexercise}
\begin{solution}
\(M_x\cong k[x,x^{-1}]\)、\(M_y\cong k[y,y^{-1}]\)，而同时逆转 \(x,y\) 会使 \(xy=0\) 变成单位为零，故 \(M_{xy}=0\)。每个 \(M\) 元素可唯一写成 \(c+xf(x)+yg(y)\)，它在两个局部化中的像为 \((c+xf(x),c+yg(y))\)。这个映射单射，且其像恰是常数项相等的多项式对。因此零次上同调为零，所有二次及更高上同调也为零；一次上同调是这份直和模掉上述像。

令 \(\delta\) 为 \((1,0)\) 的类，则 \((0,1)\) 的类等于 \(-\delta\)。全部 \(k\) 基为
\[
\delta,\quad (x^{-a},0)\ (a\ge1),\quad(0,y^{-b})\ (b\ge1).
\]
正次幂可从像中消去，负次幂不可能出现在像中，常数对则只商去对角线，故这些类线性无关且穷尽商。常数差 \(\delta\) 记录了两分支在原点处的兼容条件；它被 \(x,y\) 同时零化，但仍是一次上同调中的类，并非 \(M\) 自身的幂挠元素。
\end{solution}

\begin{optionalexercise}\label{b3:ex:22-flat-base-change}
设 \(k\) 为域，\(R=k[x]\)、\(I=(x)\)、\(B=R/(x)\)。计算 \(H_I^i(R)\otimes_RB\) 与 \(H_{IB}^i(B)\)，指出不平坦的基变换在哪个次数破坏局部上同调公式。再核对两项 Čech 复形张量后的各项，说明逐项基变换仍成立为何不足以保证上同调基变换。
\end{optionalexercise}
\begin{solution}
在 \(R\) 上，\(H_I^0(R)=0\)、\(H_I^1(R)=R_x/R\)，其余为零。乘 \(x\) 在 \(R_x/R\) 上满射，所以 \((R_x/R)\otimes_RB=(R_x/R)/x(R_x/R)=0\)。左边全部为零。\(IB=0\)，故 \(\Gamma_{IB}\) 是恒等函子，\(H_{IB}^0(B)=B\)，所有正次为零；公式在零次失败。

原 Čech 复形 \(R\to R_x\) 的张量为 \(B\to0\)，因为 \(R_x\otimes_RB=B_x=0\)。它也正是由 \(x\) 在 \(B\) 中的像构造的 Čech 复形，逐项比较并没有失败。但取上同调后得到 \(B\)，与先取原复形的零次上同调再张量所得的零模不同。不平坦张量没有保持原单射 \(R\to R_x\) 的核，这正是差别所在。
\end{solution}

\begin{optionalexercise}\label{b3:ex:22-remove-finite-torsion}
设 \((R,\mathfrak m)\) 为交换 Noether 局部环，\(M\) 为有限生成模，\(T=\Gamma_{\mathfrak m}(M)\)。证明 \(T\) 有有限长度，且 \(\Gamma_{\mathfrak m}(M/T)=0\)。若 \(M/T\ne0\)，推出它的深度至少为一。
\end{optionalexercise}
\begin{solution}
有限生成模 \(M\) 中的升链 \(0:_M\mathfrak m\subseteq0:_M\mathfrak m^2\subseteq\cdots\) 最终稳定，故 \(T\) 被某个 \(\mathfrak m^a\) 零化；其幂滤过的相邻商为有限维剩余域空间，所以长度有限。若 \(m+T\) 被 \(\mathfrak m^b\) 零化，则 \(\mathfrak m^bm\subseteq T\)，再乘 \(\mathfrak m^a\) 得零。因此 \(m\in T\)，所求商的挠子模为零。商非零时，深度与局部上同调的首次非零次数相等，而零次已消失，故深度至少一。
\end{solution}

\begin{optionalexercise}\label{b3:ex:22-dvr-dual-maps}
设 \(k\) 为域。令 \(R=k[\![t]\!]\)、\(E=k(\!(t)\!)/k[\![t]\!]\)、\(D_R=\operatorname{Hom}_R(-,E)\)，并设 \(n\ge1\) 为整数。把 \(D_R(R/(t^n))\) 写成 \(E\) 的具体子模，并求商映射 \(R/(t^{n+1})\to R/(t^n)\) 的 Matlis 对偶。
\end{optionalexercise}
\begin{solution}
同态由一的像决定，该像必须被 \(t^n\) 零化。因此
\[
D_R(R/(t^n))=(0:_Et^n)=t^{-n}R/R,
\]
其基为 \(t^{-1},\ldots,t^{-n}\) 的类。商映射的对偶是把同一个一的像看作被 \(t^{n+1}\) 零化的元素，所以是自然包含
\(t^{-n}R/R\hookrightarrow t^{-(n+1)}R/R\)。
若分别用 \(R/(t^n)\to t^{-n}R/R\)、\(a\mapsto at^{-n}+R\) 作坐标，则这个包含对应 \(a\mapsto ta\)。因此有限商的逆系统被对偶变为带乘 \(t\) 映射的正向系统。
\end{solution}

\begin{optionalexercise}\label{b3:ex:22-matlis-presentation}
设 \((R,\mathfrak m)\) 为交换完备 Noether 局部环，\(\kappa=R/\mathfrak m\)、\(E=E_R(\kappa)\)、\(D_R=\operatorname{Hom}_R(-,E)\)。有限生成的 \(R\) 模 \(M\) 有展示 \(R^a\xrightarrow{A}R^b\to M\to0\)。用矩阵 \(A\) 表示 \(D_R(M)\)。再设 \(k\) 为域，在 \(R=k[\![t]\!]\) 上计算 \(D_R(R/(t^2)\oplus R)\)。
\end{optionalexercise}
\begin{solution}
Matlis 对偶反变正合，所以
\[
0\longrightarrow D_R(M)\longrightarrow E^b\xrightarrow{A^{\mathsf T}}E^a
\]
正合，所求为转置矩阵作用的核。第二问中得到
\(D_R(R/(t^2)\oplus R)=t^{-2}R/R\oplus E\)。
第一直和项长度二，第二项 Artin 却非有限长度；因此有限生成模的对偶不能一律称作有限长度模。
\end{solution}

\begin{optionalexercise}\label{b3:ex:22-completion-in-double-dual}
设 \(k\) 为域，\(R=k[t]_{(t)}\)、\(E=k(t)/R\)、\(D_R=\operatorname{Hom}_R(-,E)\)，并取整数 \(q\ge1\)。对 \(M=R/(t^q)\oplus R\)，计算 \(D_R(D_R(M))\) 及自然求值映射的余核。解释有限长度直和项与自由直和项在非完备环上为何有不同的双对偶表现。
\end{optionalexercise}
\begin{solution}
\(E\) 是剩余域的内射包络。\Cref{b3:lem:matlis-finite-length}的有限长度求值同构把第一直和项的双对偶识别为自身，内射包络自同态的计算则把 \(D_R(D_R(R))\) 识别为 \(\widehat R=k[\![t]\!]\)。有限直和与两次对偶相容，因此
\[
D_R(D_R(M))\cong R/(t^q)\oplus k[\![t]\!].
\]
在这个识别下，求值映射在第一项上为恒等，在第二项上为完备化嵌入 \(R\hookrightarrow k[\![t]\!]\)，故余核为 \(k[\![t]\!]/R\)，非零。有限长度模已经只依赖某个有限阶商，完备化不会增加这项的数据；自由直和项则允许无限阶相容系数，所以双对偶返回其完备化。
\end{solution}

\begin{optionalexercise}\label{b3:ex:22-finite-length-local-duality}
设 \((S,\mathfrak n)\) 为 \(n\) 维交换完备 Noether 正则局部环，\(I\subsetneq S\) 为理想，\(R=S/I\) 为 \(d\) 维 CM 局部环。记 \(\mathfrak m=\mathfrak n/I\)、\(\kappa=R/\mathfrak m\)、\(D_R=\operatorname{Hom}_R(-,E_R(\kappa))\)，并取典范模 \(\omega_R=\operatorname{Ext}_S^{n-d}(R,S)\)。若 \(L\) 为有限长度 \(R\) 模，证明
\[
\operatorname{Ext}_R^j(L,\omega_R)=0\quad(j\ge0,\ j\ne d),\qquad
\operatorname{Ext}_R^d(L,\omega_R)\cong D_R(L).
\]
比较后一模与 \(L\) 的长度。
\end{optionalexercise}
\begin{solution}
\(L\) 有有限长度，故存在 \(q\) 使 \(\mathfrak m^qL=0\)。极大理想 \(\mathfrak m\) 的每个生成元 \(x_i\) 都满足 \(x_i^qL=0\)；局部化后 \(x_i\) 可逆，所以 \(L_{x_i}=0\)。因此 Čech 复形只有次数零的 \(L\)，其零次上同调为 \(L\)，正次上同调为零。局部对偶于是给出所列公式。Matlis 对偶正合且将剩余域送到同构的剩余域；对 \(L\) 的组成列逐项取对偶，因子数不变，故 \(\ell_R(D_R(L))=\ell_R(L)\)。
\end{solution}

\begin{optionalexercise}\label{b3:ex:22-quadric-canonical}
设 \(k\) 为域。在 \(S=k[\![x,y,z]\!]\) 中令 \(R=S/(xy-z^2)\)。计算其典范模并证明 \(R\) 为 Gorenstein 环；再说明它不是正则局部环。此结论是否要求 \(\operatorname{char}k\ne2\)？
\end{optionalexercise}
\begin{solution}
环境环为三维正则局部整环，\(f=xy-z^2\ne0\) 为非零因子，故商为二维 CM 环。两项消解 \(0\to S\xrightarrow{f}S\to R\to0\) 的对偶在次数一的余核为 \(R\)，所以 \(\omega_R=\operatorname{Ext}_S^1(R,S)\cong R\)，由典范模判据为 Gorenstein 环。关系属于极大理想的平方，因此商环的 \(\mathfrak m/\mathfrak m^2\) 仍为三维，而环维数为二，故不正则。所有论证对任意特征成立，没有除以二的步骤。
\end{solution}

\begin{optionalexercise}\label{b3:ex:22-artinian-pairing}
设 \(k\) 为域。令 \(A=k[u,v]/(u^2,v^3)\)，取有序基 \(1,u,v,uv,v^2,uv^2\)。写出配对 \((a,b)\mapsto\lambda(ab)\) 的矩阵，其中 \(\lambda\) 提取 \(uv^2\) 的系数。由此确定 \(A\) 的基座及 Gorenstein 性。
\end{optionalexercise}
\begin{solution}
基中元素的互补对分别是 \(1,uv^2\)、\(u,v^2\)、\(v,uv\)，因此矩阵为反对角线上全为一、其余为零的 \(6\times6\) 矩阵，行列式为 \((-1)^{15}\ne0\)。被 \(u,v\) 同时零化的元素恰为 \(k uv^2\)：乘 \(u\) 为零先去掉不含 \(u\) 的各项，再乘 \(v\) 为零只留下最高 \(v\) 次数。故基座一维，环为 Gorenstein；配对直接实现 \(A\cong\operatorname{Hom}_k(A,k)\)。
\end{solution}

\begin{optionalexercise}\label{b3:ex:22-three-generators-gorenstein}
设 \(k\) 为域。令
\[
A=k[x,y,z]/(xy,xz,yz,x^2-y^2,x^2-z^2).
\]
证明 \(A\) 是作为 \(k\) 向量空间五维的局部 \(k\) 代数，求其基座，判断其 Gorenstein 性。
\end{optionalexercise}
\begin{solution}
写 \(s=x^2=y^2=z^2\)。所有三次单项式为零：含不同变量的乘积由前三个关系为零，而 \(x^3=xy^2=0\)，其余同理。次数零、一、二分别有维数一、三、一，故基为 \(1,x,y,z,s\)。正次理想幂零，商为 \(k\)，所以它是唯一极大理想。若 \(a=a_0+a_1x+a_2y+a_3z+a_4s\) 被 \(x,y,z\) 同时零化，则先得 \(a_0=0\)，再由三个乘积分别得 \(a_1=a_2=a_3=0\)。因此基座为 \(ks\)，由基座判据为 Gorenstein 环。此计算没有根据“有几个写出的方程”判定完全交性质。
\end{solution}

\begin{optionalexercise}\label{b3:ex:22-cm-not-gorenstein-positive-dimension}
设 \(k\) 为域，令 \(R=k[\![u,v,w]\!]/(u^2,uv,v^3)\)、\(\mathfrak m=(u,v,w)R\)。计算 \(\dim R\)、\(\operatorname{depth}R\) 和 \(\dim_k\operatorname{Ext}_R^1(k,R)\)，并判断 CM 性与 Gorenstein 性。
\end{optionalexercise}
\begin{solution}
作为 \(k[\![w]\!]\) 模，\(R\) 以 \(1,u,v,v^2\) 为自由基，故 \(w\) 为非零因子。根理想为 \((u,v)\)，所以 \(\dim R=1\)，深度也为一，\(R\) 为 CM 环。商 \(B=R/wR\) 的基仍为 \(1,u,v,v^2\)，其基座为 \(ku\oplus kv^2\)：乘 \(u\) 为零迫使常数项为零，再乘 \(v\) 为零迫使一次 \(v\) 项为零。正则元素换环给出
\[
\operatorname{Ext}_R^1(k,R)\cong\operatorname{Hom}_B(k,B)
=\operatorname{Soc}(B),
\]
所以 CM 类型为二。\Cref{b3:thm:gorenstein-bass-criterion}的类型一判据于是排除 Gorenstein 性；也可由正则商保持 Gorenstein 性与 \(B\) 的二维基座得出相同结论。这里深度已经达到维数，失败来自类型而非 CM 条件。
\end{solution}

\begin{optionalexercise}\label{b3:ex:22-annihilator-duality}
设 \(k\) 为域，\(A\) 为 \(k\) 上的有限维交换局部 Gorenstein 代数，且剩余域为 \(k\)，\(J\subseteq A\) 为理想。证明
\[
\dim_k(0:_AJ)=\dim_k(A/J),\qquad
0:_A(0:_AJ)=J.
\]
\end{optionalexercise}
\begin{solution}
取正文中的非退化乘法配对 \(\lambda(ab)\)。\(J\) 的正交补恰是 \((0:_AJ)\)：一个方向直接成立；反之，若 \(\lambda(aJ)=0\) 但 \(aJ\ne0\)，非零理想 \(aJ\) 与一维基座相交，而 \(\lambda\) 在基座上非零，矛盾。非退化有限维配对给出正交补维数等于总维数减去 \(\dim_kJ\)，得到第一式；再取正交补恢复原子空间，得到第二式。理想条件用于使 \(aJ\) 仍为理想，不能把零化子结论照搬到任意线性子空间。
\end{solution}

\begin{optionalexercise}\label{b3:ex:22-longer-embedded-component}
设 \(k\) 为域。在 \(S=k[\![x,y]\!]\) 中令 \(R=S/(x^2,xy^2)\)、\(\mathfrak m=(x,y)R\)。求 \(H_{\mathfrak m}^0(R)\)、\(H_{\mathfrak m}^1(R)\)，以及 \(\operatorname{Ext}_S^1(R,S)\)、\(\operatorname{Ext}_S^2(R,S)\)。比较两组计算。
\end{optionalexercise}
\begin{solution}
每个元素唯一写成 \(a(y)+x(b_0+b_1y)\)。局部化于 \(y\) 后 \(x=0\)，所以
\[
H_{\mathfrak m}^0(R)=kx\oplus kxy,\qquad
H_{\mathfrak m}^1(R)=k(\!(y)\!)/k[\![y]\!].
\]
这里 \(\sqrt{(y)}=\mathfrak m\)，故两项 Čech 复形足够。另一方面，两生成元的全部关系由 \((-y^2,x)\) 生成，得到长度二的环境环自由消解。对偶复形中，\(-y^2a+xb=0\) 等价于 \((a,b)=(xc,y^2c)\)，前一像是此核的 \(x\) 倍，故
\[
\operatorname{Ext}_S^1(R,S)=S/(x),\qquad
\operatorname{Ext}_S^2(R,S)=S/(x,y^2).
\]
第二模长度二，与零次局部上同调的 Matlis 对偶长度相同；第一模对应一次局部上同调。比正文长度一的嵌入部分，这个例子还记录了非零作用 \(y\cdot x=xy\)，不只记录一个维数。
\end{solution}

\begin{optionalexercise}\label{b3:ex:22-cm-module-environment-ext}
设 \((S,\mathfrak n)\) 为 \(n\) 维交换完备 Noether 正则局部环，\(M\ne0\) 为有限生成的 \(S\) 模，\(r=\dim_SM\)。证明 \(M\) 为 CM 模当且仅当 \(\operatorname{Ext}_S^j(M,S)\) 仅在 \(j=n-r\) 非零。
\end{optionalexercise}
\begin{solution}
正则局部对偶把 \(\operatorname{Ext}_S^j(M,S)\) 识别为 \(D_S(H_{\mathfrak n}^{n-j}(M))\)。Matlis 对偶是忠实的：非零模的非零元素可由映入内射包络的同态检测，所以一个模为零当且仅当其对偶为零。因此 Ext 集中在 \(n-r\) 等价于局部上同调集中在 \(r\)。由深度首次非零和维数以上消失的两个定理，这又恰好等价于 \(\operatorname{depth}_SM=r\)。在两种情形中所指的唯一一项确实非零，而不只是说其他项消失。
\end{solution}
